\ReportCover{基于 Schema 的 Kernel 表示与生成}{TRE AI系统工程}{邵瀚雍}{岩泉｜北园}{2026 年 11 月 26 日}
\ReportContents
\SourceSeparator
\clearpage
\section{自动化算子生成的研究意义}\label{sec:motivation}
高性能算子（Kernel）连接上层的深度学习算法和底层的异构硬件平台。我们做自动化算子生成的意义，看重的不只是得到一个局部的代码生成工具。这件事同时对应两方面的需要：面向当下，MaaS 大规模推理需要更多、更快的算子；面向未来，它非常适合作为研究 AI 自我改进闭环的起点。


\begin{DiagramBlock}
\IfFileExists{diagrams/diagram-01.tex}{%
\begin{adjustbox}{max width=\linewidth,max totalheight=.80\textheight,center}
\begingroup%
\definecolor{ksBlue}{HTML}{3478B8}%
\definecolor{ksTeal}{HTML}{27877F}%
\definecolor{ksAmber}{HTML}{C58B28}%
\definecolor{ksInk}{HTML}{243447}%
\definecolor{ksPaleBlue}{HTML}{F0F6FC}%
\definecolor{ksPaleTeal}{HTML}{EFF8F5}%
\begin{tikzpicture}[
  x=1mm, y=1mm,
  card/.style={draw, fill=white, rounded corners=2mm, line width=.7pt,
    text=ksInk, align=flush center, font=\fontsize{9.5}{12.5}\selectfont,
    inner xsep=2mm, inner ysep=1.7mm, outer sep=.35pt},
  future/.style={card, draw=ksBlue, minimum width=69mm},
  current/.style={card, draw=ksTeal},
  frame/.style={draw, rounded corners=3mm, line width=.7pt, inner sep=3mm,
    outer sep=.35pt},
  heading/.style={text=ksInk, font=\bfseries\fontsize{10}{12.5}\selectfont,
    inner sep=0pt, outer sep=0pt},
  flow/.style={-{Stealth[length=2.25mm,width=1.65mm]}, line width=.7pt,
    rounded corners=2mm}
]
  \node[heading] (future-title) at (83,0) {面向未来：可验证的自动化闭环};
  \node[heading] (now-title) at (0,0) {面向当下：MaaS 推理的算子需求};

  \node[future] (f1) at (83,-14) {模型能力提升的新来源\\
    RL 合成数据 · 系统自我迭代 (RSI)};
  \node[future] (f2) at (83,-32) {前提：边界清楚、\\
    结果能客观验证的好问题};
  \node[future] (f3) at (83,-53) {算子开发天然满足这一前提\\
    正确性可判定（编译 + 数值对齐）\\
    性能反馈连续（延迟 / Roofline）\\
    软硬件协同的组合优化，有足够难度};

  \node[current, minimum width=38.5mm, minimum height=23mm] (n1) at (-20.5,-20)
    {模型架构快速演进\\
    Dense / MoE / MLA\\
    线性与稀疏注意力\\
    FP8 / MXFP8 / NVFP4};
  \node[current, minimum width=38.5mm, minimum height=23mm] (n2) at (20.5,-20)
    {异构推理卡多代并存\\
    Hopper / Blackwell\\
    Ascend 等国产卡\\
    HIP 等其他平台};
  \node[current, minimum width=75mm] (n3) at (0,-43)
    {算子需求呈 M × N × K 组合爆炸\\
    HPC 专家少，手工调优周期长};
  \node[current, minimum width=75mm] (n4) at (0,-59)
    {算子效率直接影响\\
    单 Token 成本 · 算力利用率 · 服务延迟};

  \coordinate (future-base) at (83,-65);
  \coordinate (now-base) at (0,-65);
  \begin{scope}[on background layer]
    \node[frame, draw=ksBlue, fill=ksPaleBlue,
      fit=(future-title)(f1)(f2)(f3)(future-base)] (future-group) {};
    \node[frame, draw=ksTeal, fill=ksPaleTeal,
      fit=(now-title)(n1)(n2)(n3)(n4)(now-base)] (now-group) {};
  \end{scope}

  \node[card, draw=ksAmber, line width=.85pt, minimum width=96mm] (core)
    at (41.5,-87) {目标：构建可验证、能收敛、可理解的\\
    自动化算子生成闭环，并把方法用到更多任务上};
  \coordinate (future-turn) at (62,-76);
  \coordinate (now-turn) at (21,-76);

  \draw[flow, ksBlue] (f1) -- (f2);
  \draw[flow, ksBlue] (f2) -- (f3);
  \draw[flow, ksTeal] (n1.south) -- (n1.south |- n3.north);
  \draw[flow, ksTeal] (n2.south) -- (n2.south |- n3.north);
  \draw[flow, ksTeal] (n3) -- (n4);
  \draw[flow, ksBlue] (future-group.south) |- (future-turn) --
    (future-turn |- core.north);
  \draw[flow, ksTeal] (now-group.south) |- (now-turn) --
    (now-turn |- core.north);
\end{tikzpicture}%
\endgroup%
%
\end{adjustbox}%
}{%
\MissingDiagram{01}%
}
\end{DiagramBlock}
\subsection{面向当下：MaaS 推理的算子需求}\label{sec:maas-demand}
首先我们聚焦于现实的需求。在大规模 MaaS（Model-as-a-Service）推理中，算子需求是下面三类因素的笛卡尔积，每多一种模型结构、一款卡或一种负载，要写的算子就成倍增加：


\begin{itemize}
\item
\textbf{模型结构与精度}：MHA / GQA、MLA、线性注意力（GDN / KDA）、稀疏注意力、Fused MoE 等结构，再叠加 BF16、FP8、MXFP8、NVFP4 等量化格式；


\item
\textbf{硬件代际}：要用好 Blackwell（B200 / B300）的 \InlineCode{tcgen05}、TMEM、2-CTA 等新原语，也要适配国产推理卡和不同代际的集群；


\item
\textbf{负载形态}：Prefill、Decode、MTP 等阶段以及不同序列长度，对吞吐和延迟的要求各不相同。


\end{itemize}

算子效率直接影响单 Token 成本、算力利用率和 TTFT / TPOT 延迟。同时懂数学推导和微架构的 HPC 专家很少，一个算子在一款卡上手写调优要几周到几个月，跟不上模型与硬件组合增长的速度。


\subsection{面向未来：算子开发是可验证的好问题}\label{sec:verifiable-loop}
除了眼前的需求，这件事还有更长远的意义。预训练语料带来的增益正在放缓，模型能力的提升越来越依赖强化学习（RL）合成数据、长程工具交互（Agentic Harness）和系统的自迭代、自进化（Recursive Self-Improvement）。这类自动化循环要转起来，首先需要一个\textbf{边界清楚、结果能客观验证的好问题（Good Problem \& Verifiable Loop）}：先在这个问题上把闭环建起来，再推广到各类任务。在众多工程和科研任务中，算子开发是少有的满足这个条件的问题，适合作为研究自动化闭环的起点，有三个关键理由：


\begin{enumerate}[start=1]
\item
\textbf{结果可验证}：能否编译、数值是否与参考实现（Reference）一致，都有确定的判据；实测延迟和 Roofline 利用率又给出来自真实硬件的连续反馈，每次尝试都能拿到一个分数。


\item
\textbf{问题有难度}：写好算子要在寄存器、片上存储（SMEM / TMEM）和带宽的约束下，同时决定数学等价改写、并行切分、多级存储复用，以及用异步流水掩盖访存延迟。这是一个组合优化问题，靠语法翻译解决不了。


\item
\textbf{方法可迁移}：如果能在算子这个软硬件协同问题上打磨出一套高效、能收敛、能积累经验的自动化 Loop，其中问题表示、分层搜索、门禁验证和经验蒸馏等机制有望用到其他系统设计和自动化科研任务上。


\end{enumerate}


\textbf{向近看，自动化算子生成是降低 MaaS 推理成本的必经之路；向远看，它是探索 AI 自我改进闭环的理想起点。}


\SourceSeparator
\clearpage
\section{现有的研究现状和局限性}\label{sec:landscape}
正因为自动化算子生成兼有研究价值和落地价值，近两年学术界、大模型公司和芯片厂商都在做 Kernel Agent，工作覆盖数据集、评测基准、专用模型训练和多智能体搜索框架。


\subsection{三条主流技术路线}\label{sec:three-routes}
按模型能力从哪里来、搜索如何组织，现有工作大致分为三条路线：


\begin{RouteCard}
\CardField{技术路线}{\textbf{路线一：垂域后训练型} \emph{(LLM4Kernel / Agentic RL)}}
\CardField{代表工作}{\textbf{CUDA Agent} \emph{(字节 Seed × 清华, 2026)}}
\CardField{实现方法}{合成多样的训练题，在能编译、测试和 Profiling 的真实 CUDA 环境里做多轮 Agentic RL，把排障和改写代码的能力训练进模型权重}
\CardField{优势}{端到端调用工具和单点优化的能力明显提升；分档奖励和随机输入减少了模型钻评测漏洞（Reward Hacking）的情况}
\CardField{局限}{需要昂贵的可执行训练环境和题目；能力固化在权重里，绑定特定硬件和语言分布，难以迁移到新架构}
\end{RouteCard}
\begin{RouteCard}
\CardField{技术路线}{\textbf{路线二：推理期搜索与执行反馈} \emph{(Agent4Kernel)}}
\CardField{代表工作}{\textbf{AlphaEvolve} \emph{(DeepMind)} · \textbf{AVO} \emph{(NVIDIA, 2026)} · \textbf{KDA} \emph{(MIT HAN Lab, 2026)} · \textbf{AKA} \emph{(内部 Atrex-Kernel-Agent)}}
\CardField{实现方法}{冻结强基座模型，结合进化搜索 / 候选池、硬件知识库（Wiki）、Profiler 反馈（如 NCU）、独立校验和按真实负载 Roofline 加权的评分，做长程迭代}
\CardField{优势}{在多个单算子任务上提速明显：AVO 在 B200 上从强先验出发连续演化 7 天，接近手写库，在部分算子上超过；KDA 显示加入 Wiki 和 NCU 反馈后加速比大幅提高}
\CardField{局限}{直接在底层代码上做变异搜索，每轮试错周期长，Token 和 GPU 开销大；收敛随机性强，产物容易变成黑盒，分支越来越多}
\end{RouteCard}
\begin{RouteCard}
\CardField{技术路线}{\textbf{路线三：高层意图 / 调度中间层} \emph{(Sketch / IR 分层)}}
\CardField{代表工作}{\textbf{K-Search} \emph{(UC Berkeley, 2026)} · \textbf{AKG Kernel Agent} \emph{(华为)} · \textbf{CAKE} \emph{(NVIDIA × CMU, 2026)} · \mbox{\textbf{ForgeKernel}} \emph{(面壁智能, 2026)}}
\CardField{实现方法}{把高层策略和底层代码分开：K-Search 维护一棵自然语言优化意图树；AKG 生成跨 DSL 的 Unified Sketch；CAKE 定义编译器级调度 IR，再确定性地生成代码；ForgeKernel 由架构 Agent 制定里程碑计划，编码 Agent 按计划实现}
\CardField{优势}{一次编译失败不会丢掉高层优化意图；ForgeKernel 在 Hopper / Blackwell 的 29 个配置中有 25 个持平或超过参考实现，初步说明策略与实现分离是可行的}
\CardField{局限}{自然语言意图和粗粒度 Sketch 难以精确约束片上资源、同步和数学改写；CAKE 一类 IR 绑定特定算子域和专用后端，难以跨语言通用，也没有给具体实现留白；ForgeKernel 依赖已有的高性能参考实现，缺少参考时优化容易停滞}
\end{RouteCard}
\subsection{共同局限：以算子coding为中心的编程视角}\label{sec:common-limits}
这些工作说明，LLM 和 Agent 闭环可以生成可用的算子。但把视角从中小规模标准算子的刷榜，拉回到复杂生产算子（如 Blackwell 上的 FA4、Fused MoE、MLA、GDN）从研发到上线的完整流程，现有 Kernel Agent 无论换用哪种搜索技巧，都会碰到完全相同的问题。


而关键在于表示方法：\textbf{现有方案大多把未经压缩的底层代码同时当作知识的存放形式、搜索空间和交付产物}。按算子生成的流程看，这种做法在“数据与表示”、“生成与迭代”、“交付与维护”三个环节各自带来问题，三者环环相扣；我们在下一章展开。


\SourceSeparator
\clearpage
\section{自动化算子生成面临的三个挑战}\label{sec:challenges}
刚刚提到，现有方案的问题分别出现在算子生成流程的数据与表示、生成与迭代、交付与维护三个环节，对应的依次限制了自动设计的能力上限、迭代的收敛效率，以及产物能否真正进入生产环境。我们将在本章详细展开。


\begin{DiagramBlock}
\IfFileExists{diagrams/diagram-02.tex}{%
\begin{adjustbox}{max width=\linewidth,max totalheight=.80\textheight,center}
\begingroup%
\definecolor{ksBlue}{HTML}{3478B8}%
\definecolor{ksTeal}{HTML}{27877F}%
\definecolor{ksAmber}{HTML}{C58B28}%
\definecolor{ksInk}{HTML}{243447}%
\definecolor{ksPaleBlue}{HTML}{F0F6FC}%
\definecolor{ksPaleTeal}{HTML}{EFF8F5}%
\definecolor{ksPaleAmber}{HTML}{FFF8E8}%
\begin{tikzpicture}[
  x=1mm, y=1mm,
  card/.style={draw, fill=white, rounded corners=2mm, line width=.7pt,
    text=ksInk, align=flush center, font=\fontsize{9.5}{12.5}\selectfont,
    text width=69mm, minimum width=73mm, minimum height=14mm, inner xsep=2mm, inner ysep=1.5mm,
    outer sep=.35pt},
  challenge/.style={card, font=\bfseries\fontsize{9.5}{12.5}\selectfont},
  frame/.style={draw, rounded corners=3mm, line width=.7pt, inner sep=3mm,
    outer sep=.35pt},
  heading/.style={text=ksInk, font=\bfseries\fontsize{10}{12.5}\selectfont,
    inner sep=0pt, outer sep=0pt},
  link/.style={line width=.65pt},
  flow/.style={-{Stealth[length=2.6mm,width=1.9mm]}, ksInk, line width=.9pt},
  edge text/.style={text=ksInk, font=\fontsize{9}{12.5}\selectfont,
    inner sep=1mm}
]
  \node[heading] (s1-title) at (40,0) {环节一：数据与表示};
  \node[challenge, draw=ksBlue] (c1-title) at (0,-12)
    {挑战一：硬件感知不足};
  \node[card, draw=ksBlue] (c1-a) at (80,-12)
    {语料偏置：以 NVIDIA 旧架构为主\\
    Blackwell 与国产卡语料稀缺};
  \node[card, draw=ksBlue] (c1-b) at (80,-29)
    {语言杂（语言是思维的边界）：\\
    CUDA / Triton / TileLang / Ascend 等};
  \node[card, draw=ksBlue] (c1-c) at (0,-29)
    {内容杂（源码未经压缩与蒸馏）：\\
    Host / IO 样板淹没设计决策};

  \node[heading] (s2-title) at (40,-51) {环节二：生成与迭代};
  \node[challenge, draw=ksTeal] (c2-title) at (0,-63)
    {挑战二：迭代收敛困难};
  \node[card, draw=ksTeal] (c2-a) at (80,-63)
    {随机性高、一致性低：\\
    同条件重复运行，轨迹明显分化};
  \node[card, draw=ksTeal] (c2-b) at (80,-80)
    {设计与实现强耦合：\\
    好设计因实现瑕疵被“误杀”};
  \node[card, draw=ksTeal] (c2-c) at (0,-80)
    {避难就易走弯路：为尽快跑通\\
    退回 mma.sync 等旧指令，平台期长};

  \node[heading] (s3-title) at (40,-102) {环节三：交付与维护};
  \node[challenge, draw=ksAmber] (c3-title) at (0,-114)
    {挑战三：维护实用麻烦};
  \node[card, draw=ksAmber] (c3-a) at (80,-114)
    {Shape 特化分支膨胀：\\
    以 if-else 针对特定负载“刷分”，缺乏泛化};
  \node[card, draw=ksAmber] (c3-b) at (80,-131)
    {隐式缺陷：共享内存竞争、越界、死锁\\
    测试集随之膨胀，迭代变慢};
  \node[card, draw=ksAmber] (c3-c) at (0,-131)
    {黑盒代码难以审阅：\\
    设计取舍埋在底层代码中，上线风险高};

  \begin{scope}[on background layer]
    \node[frame, draw=ksBlue, fill=ksPaleBlue,
      fit=(s1-title)(c1-title)(c1-a)(c1-b)(c1-c)] (s1) {};
    \node[frame, draw=ksTeal, fill=ksPaleTeal,
      fit=(s2-title)(c2-title)(c2-a)(c2-b)(c2-c)] (s2) {};
    \node[frame, draw=ksAmber, fill=ksPaleAmber,
      fit=(s3-title)(c3-title)(c3-a)(c3-b)(c3-c)] (s3) {};
  \end{scope}

  \draw[link, ksBlue] (c1-title) -- (c1-a);
  \draw[link, ksBlue] (c1-a) -- (c1-b);
  \draw[link, ksBlue] (c1-b) -- (c1-c);
  \draw[link, ksTeal] (c2-title) -- (c2-a);
  \draw[link, ksTeal] (c2-a) -- (c2-b);
  \draw[link, ksTeal] (c2-b) -- (c2-c);
  \draw[link, ksAmber] (c3-title) -- (c3-a);
  \draw[link, ksAmber] (c3-a) -- (c3-b);
  \draw[link, ksAmber] (c3-b) -- (c3-c);
  \draw[flow] (s1.south) --
    node[edge text, right=1mm] {带着不完整的先验进入搜索} (s2.north);
  \draw[flow] (s2.south) --
    node[edge text, right=1mm] {产出高风险的黑盒代码} (s3.north);
\end{tikzpicture}%
\endgroup%
%
\end{adjustbox}%
}{%
\MissingDiagram{02}%
}
\end{DiagramBlock}

\newpage

\subsection{数据与表示：模型懂编程，不懂硬件}\label{sec:challenge-data}
现有基座模型和编码 Agent 普遍“重编程、轻硬件”：通用编程与代码改写能力较强，但对硬件的理解和软硬件协同设计（Algorithm-Hardware Co-design）能力较弱。它们能写出语法正确的核函数，却不清楚目标硬件的吞吐、存储层级和延迟掩盖机制，也说不清为什么要采用某种切分和流水结构。原因主要在语料和表达方式上，有以下三点。


\begin{enumerate}[start=1]
\item
\textbf{训练语料架构代际混杂}


语料集中在 NVIDIA 的历史架构（如 Ampere、Hopper），不同代际的实现混在一起，也没有标注各代硬件的能力边界。面对新微架构（如 Blackwell 的 TMEM 与异步指令）或昇腾 Ascend、AMD HIP 等平台时，模型缺少目标卡的可靠先验，常常照搬旧卡的模板。国产模型和跨架构迁移场景受影响最大。


\item
\textbf{编程语言细节碎片化}


算子开发语言很多，包括 \InlineCode{CUDA C++}、\InlineCode{PTX}、\InlineCode{CuTe /\allowbreak{} CuTeDSL}、\InlineCode{Triton}、\InlineCode{TileLang}、\InlineCode{Ascend C}、\InlineCode{HIP} 等，抽象层级各不相同，规范标准也完全不同。

\textbf{语言是思维的边界}：Agent 在某一种面向特定硬件抽象的语言里思考时，设计空间就被框定在这种语言的惯用写法之内，很难看到各语言之间共通的数学改写和调度结构。Megakernel 这类打破单算子边界、需要重新划定切分与融合范围的场景，受这一限制最明显。


\item
\textbf{工程实现细节耦合设计决策}


现有算子源码大多没有经过整理，信息密度很低。一份工业级算子动辄数千行（如 FlashAttention-4 前向核心 CUDA 代码超过 2300 行），大部分篇幅是 host 衔接层、张量 IO 校验、指针步长、模板元编程与语法样板；\textbf{决定性能上限的设计决策（数学改写、工作块映射、片上缓冲复用、异步流水编排）只有百余行，隐藏在这些代码中间}。

\textbf{压缩即智能}：大模型的能力来自对数据中共性规律的压缩。源码中与设计无关的工程代码占了大头，模型能从中提炼出的软硬件协同规律就很少。


\end{enumerate}

总的来说，现有算子源代码信息密度低，导致同一类设计散落在多种语言里，每份源码中讲设计的部分又很少。于是，训练时的语料和推理时查阅的 Wiki 都处在未压缩、未蒸馏的状态。表达设计的语言不统一、不凝练，造成了数据问题，也限制了模型思考设计的方式，“重编程、轻硬件”由此成为自动设计的能力上限。

\newpage

\subsection{生成与迭代：难以收敛，也不知何时停止}\label{sec:challenge-iteration}
在生成与迭代环节，问题是相同配置多次运行的结果不一致，容易停在平台期，而且无法判断成本与收益之间的帕累托前沿（Pareto Frontier），也就不知道该在哪里停下。


每一轮优化都要消耗实际迭代时间、GPU 机时和大量 Token，这些都是需要减少的成本。我们希望在成本可接受的位置停下，同时确信此时的解已接近该算子在目标硬件上的性能极限。这需要回答两个问题：\textbf{当前的solution离硬件极限SOL还有多远？什么时候应该停止？}前者是预测问题，后者是收敛问题。在现有的端到端代码迭代中，最优路径难以稳定复现，这两个问题都无从回答。


我们实测的两组数据可以说明这一点。在 B300 上用完全相同的配置重复运行 AKA Loop，前 10 个 episode 内，各次运行最终接受的延迟最慢与最快相差 1.44×～1.67×，单轮探索累计消耗数亿输入 Token（含缓存）。在 \InlineCode{GQA} 任务中，Opus 在第 97 个版本达到 13.34 µs；Qwen 经过很长的平台期，到第 1224 个版本才达到最佳的 13.83 µs，跑到 \InlineCode{v1367} 仍为 13.91 µs。系统无法判断这是已经接近硬件上限，还是停在了局部次优解。


\OriginalImage{assets/gdn_prefill_trace.png}{GDN Prefill 从零迭代的多路收敛轨迹}{\emph{图：L20D（SM103）上} \InlineCode{GDN Prefill} \emph{多路独立重复运行的收敛轨迹。橙色为基线组（}\InlineCode{Control r1–r3}\emph{）：初始条件相同、无外部参考，各 Run 的路径明显分化，其中} \InlineCode{Control r2} \emph{长期停在} 365.9 µs \emph{附近，与同组其他路径相差数倍。蓝色为第\ref{sec:schema}章引入 Kernel Schema 后的对照组。}}
\begin{DiagramBlock}
\IfFileExists{diagrams/diagram-03.tex}{%
\begin{adjustbox}{max width=\linewidth,max totalheight=.80\textheight,center}
\begingroup%
\definecolor{ksBlue}{HTML}{3478B8}%
\definecolor{ksAmber}{HTML}{C58B28}%
\definecolor{ksRed}{HTML}{A44C46}%
\definecolor{ksInk}{HTML}{243447}%
\definecolor{ksPaleAmber}{HTML}{FFF8E8}%
\begin{tikzpicture}[
  x=1mm, y=1mm,
  card/.style={draw, fill=white, rounded corners=2mm, line width=.7pt,
    text=ksInk, align=flush center, font=\fontsize{9.5}{12.5}\selectfont,
    inner xsep=2mm, inner ysep=1.7mm, outer sep=.35pt},
  branch/.style={card, draw=ksAmber, minimum width=50mm, minimum height=18mm},
  flow/.style={-{Stealth[length=2.25mm,width=1.65mm]}, draw=ksInk, line width=.7pt,
    rounded corners=2mm},
  edge text/.style={text=ksInk, fill=white, font=\fontsize{9}{12.5}\selectfont,
    inner xsep=1mm, inner ysep=.6mm}
]
  \node[card, draw=ksBlue, minimum width=86mm] (idea) at (0,0)
    {Agent 提出新的流水线设计\\
    （如调整 QK / PV 计算顺序）};
  \node[card, draw=ksBlue, minimum width=86mm] (rewrite) at (0,-17)
    {需同步重写上千行底层代码\\
    （索引 · 布局 · 屏障同步）};
  \node[card, diamond, rounded corners=0pt, aspect=2.65, draw=ksAmber,
    fill=ksPaleAmber, inner xsep=1mm, inner ysep=.9mm] (noise) at (0,-43)
    {实现缺陷与测量扰动\\
    编译错误 / 未调优 / 测量波动};

  \node[branch] (false-negative) at (-56,-86)
    {好设计被“误杀”\\
    (False Negative)\\
    误以为设计方向不可行};
  \node[branch] (detour) at (0,-86)
    {退回熟悉但低效的旧路线\\
    如在 Blackwell 上用\\
    mma.sync 替代 tcgen05};
  \node[branch] (diverge) at (56,-86)
    {6 路并行 Agent 路径分化\\
    跨 Run 方差大\\
    收敛终点难预测};
  \node[card, draw=ksRed, line width=.85pt, minimum width=95mm] (trap)
    at (0,-114) {停滞于局部次优平台期：耗时长，Token 开销大};

  \draw[flow] (idea) -- (rewrite);
  \draw[flow] (rewrite) -- (noise);
  \draw[flow, draw=ksAmber] (noise.215) |- (-56,-68) -- (false-negative.north);
  \draw[flow, draw=ksAmber] (noise.south) -- (detour.north);
  \draw[flow, draw=ksAmber] (noise.325) |- (56,-68) -- (diverge.north);
  \node[edge text, anchor=south] at (-42,-67.5) {实测变慢或未完成};
  \node[edge text, anchor=west] at (1,-71) {为尽快获得正反馈而保守退回};
  \node[edge text, anchor=south] at (40,-67.5) {微小差异被放大（蝴蝶效应）};
  \draw[flow, draw=ksRed] (false-negative.south) |- (-32,-105) --
    (-32,-105 |- trap.north);
  \draw[flow, draw=ksRed] (detour) -- (trap);
  \draw[flow, draw=ksRed] (diverge.south) |- (32,-105) -- (32,-105 |- trap.north);
\end{tikzpicture}%
\endgroup%
%
\end{adjustbox}%
}{%
\MissingDiagram{03}%
}
\end{DiagramBlock}
我们追踪了多路并行 Agent（如 6 个 Opus Agent 并行优化同一算子）的完整思维链和版本演化，发现主要原因是\textbf{高层设计决策与底层代码实现耦合在一起}：改一处设计就要改大量代码，而代码里的任何问题又会反过来影响 Agent 对设计的判断。具体有三种表现。


\begin{enumerate}[start=1]
\item
\textbf{微小扰动放大为路径分叉}：即使起点完全相同，一次编译报错、寄存器溢出或测量波动，也会经推理链层层放大，产生“蝴蝶效应”，让不同 Agent 走上完全不同的技术路线。


\item
\textbf{好设计因实现瑕疵被“误杀”（False Negative）}：结构性的新设计（如调整 QK 与 PV 的计算顺序、重构多级异步流水）往往要重写成百上千行代码。其中只要有局部索引错误、布局未对齐或参数未调优，实测就会变慢甚至超时。Agent 收到的负反馈分不清是“实现没写好”还是“设计不好”，于是放弃了一条本来正确且有潜力的设计。早期 FA4 测试中，一个采用 \InlineCode{tcgen05} 和新流水结构的正确方案，就因为单轮没有写对而被门禁淘汰。


\item
\textbf{为求跑通退回低效弯路}：重写成本高，Agent 倾向于先通过编译和正确性检查。例如在 Blackwell 上绕开原生的 \InlineCode{tcgen05.\allowbreak{}mma} + TMEM 异步方案，退回语料里最熟悉的 \InlineCode{mma.\allowbreak{}sync} 并反复做局部修补，搜索预算消耗在提升空间有限的分支上。


\end{enumerate}
\SourceSeparator
\subsection{交付与维护：产物是黑盒，难审难维护}\label{sec:challenge-delivery}
第三个问题在于交付的环节：产物能否让人看懂、能否可靠可控，决定了它能否进入生产。现状是 Agent 为特定输入写了大量分支特化，代码随之膨胀，内部还藏有非常深的并发和边界错误；人类专家难以界定它的行为，也难以读懂和修改，长期维护成本高。


我们在自建的 44 个算子主库（20 个 AKA 迭代解 + 24 个参考算子实现）上做了审计。对其中首批 18 个算子（共 48 份设计）做静态提取和独立的正确性复审后，发现了三类核心问题。


\begin{enumerate}[start=1]
\item
\textbf{靠 Shape 特化“刷分”，代码膨胀}


Agent 倾向于堆砌 \InlineCode{if-\allowbreak{}else} 分支和针对特定 Shape 的派生核函数，以弥补单一设计在不同输入上的性能不足。例如 \InlineCode{fmha\_\allowbreak{}mxfp8\_\allowbreak{}qkv} 拆出了数十个计划分支，仅 D128 split 一族就有“扫描 × 搬运 × 写回 × 归并”24 种组合；\InlineCode{nvfp4\_\allowbreak{}fused\_\allowbreak{}moe} 最初有 50 种 FC1 × FC2 命名配对。交付代码常达数千行，人很难看出它为什么更快。实际上，这类加速往往来自对少数评测 Shape 的特化，在能推广到其他 Shape 的设计上并没有新的突破。


\item
\textbf{并发与边界等缺陷难以感知，测试集越堆越慢}


能通过有限评测用例的代码，仍可能含有不易察觉的错误。首批 18 个算子中的典型问题有两类。一是共享内存竞争与越界：\InlineCode{fmha\_\allowbreak{}mxfp8\_\allowbreak{}qkv} 的空 split 分支中，零有效块路径仍会写入零部分和，但复用共享缓冲前缺少跨 warp 同步（后置的 CTA 屏障无法弥补）；空分支入口还存在读取 \InlineCode{MState[cLo-\allowbreak{}1]} 的越界风险。二是零工作量下的同步失配：\InlineCode{paged\_\allowbreak{}indexed\_\allowbreak{}latent\_\allowbreak{}fp8} 在 \InlineCode{active\_\allowbreak{}topk=0} 时，生产者的通知执行 0 次，而消费者仍无条件等待首个通知，存在死锁风险。另有部分算子在测试未覆盖的分支上依赖慢速 fallback。

要拦住这类缺陷，只能不断扩充测试集、增加极端用例，结果是单次评估越来越慢，迭代也被拖慢，而测试没有覆盖的边界上仍可能有缺陷。


\item
\textbf{设计意图丢失，专家难以介入和理解算子}


人类专家擅长做高层取舍，也有设计品味（Taste），例如知道为什么选 3 级流水（2 级掩盖不住 TMA 延迟，4 级超出 SMEM 预算），或者延迟重缩放阈值（如 $\tau=8$）的精度边界在哪里。端到端生成的代码把这些取舍埋在语法细节里，算子成了黑盒，专家很难以 Human-in-the-Loop 的方式参与。\textbf{要上线，需要大量人力逐行审计；直接上线，又可能因隐式竞争和边界行为引发线上故障。}这限制了自动生成的算子在生产中大规模使用。


\end{enumerate}
\SourceSeparator
\clearpage
\section{从“工程编码”回归“算子设计”的第一性原理}\label{sec:first-principles}
提出方案之前，先看一个基本问题：\textbf{一个优秀的 HPC 工程师是怎样开发算子的？编码 Agent 和他们差在哪里？}


\subsection{人类专家先做设计，编码 Agent 直接写代码}\label{sec:design-first}
对经验丰富的算子工程师来说，写 CUDA 或 CuTe 代码只是最后一步的工程翻译。\textbf{主要的工作是在给定硬件资源和数学目标构成的空间里，搜索最优的软硬件协同设计（Design）。}


以 Blackwell（B200 / B300）上的 FlashAttention-4 前向算子为例，专家在写第一行代码之前，已经完成了两步工作。


第一步是定量定位瓶颈。依据指令延迟和 Microbench 实测吞吐，核算主循环的资源账本：单块 Tensor Core UMMA（\InlineCode{QK + PV}）约 512 cycles，TMA 访存约 507 cycles，而 SFU 上的 \InlineCode{exp2} 约 1248 cycles，因此 SFU 吞吐是第一瓶颈。


第二步是围绕瓶颈安排掩盖与配平。用 Warp-Specialization 划分专职角色，让两组 Softmax Warpgroup 交错执行以分摊 SFU 压力，规划 SMEM 与 TMEM（Tensor Memory）的槽位容量和生命周期，并安排 TMA 搬运、UMMA 与在线重缩放之间的流水重叠。


这些推理决定了算子的性能，但都不涉及编程语言语法、Host 封装或指针计算，属于设计（Design）层面的决策。


\begin{DiagramBlock}
\IfFileExists{diagrams/diagram-04.tex}{%
\begin{adjustbox}{max width=\linewidth,max totalheight=.80\textheight,center}
\begingroup%
\definecolor{ksTeal}{HTML}{27877F}%
\definecolor{ksAmber}{HTML}{C58B28}%
\definecolor{ksInk}{HTML}{243447}%
\definecolor{ksPaleTeal}{HTML}{EFF8F5}%
\definecolor{ksPaleAmber}{HTML}{FFF8E8}%
\begin{tikzpicture}[
  x=1mm, y=1mm,
  card/.style={draw, fill=white, rounded corners=2mm, line width=.7pt,
    text=ksInk, align=flush center, font=\fontsize{9.5}{12.5}\selectfont,
    minimum width=60mm, minimum height=14mm, inner xsep=2mm, inner ysep=1.7mm,
    outer sep=.35pt},
  frame/.style={draw, rounded corners=3mm, line width=.7pt, inner sep=3mm,
    outer sep=.35pt},
  heading/.style={text=ksInk, font=\bfseries\fontsize{10}{12.5}\selectfont,
    inner sep=0pt, outer sep=0pt},
  flow/.style={-{Stealth[length=2.25mm,width=1.65mm]}, line width=.7pt},
  edge text/.style={text=ksInk, font=\fontsize{9}{12.5}\selectfont, inner sep=1mm}
]
  \node[heading] (human-title) at (34,0)
    {人类 HPC 专家：把算子视为 Design 搜索问题};
  \node[card, draw=ksTeal] (h1) at (0,-13)
    {Microbench 测算\\
    各硬件单元的延迟与吞吐};
  \node[card, draw=ksTeal] (h2) at (68,-13)
    {定位主要瓶颈\\
    如 B200 FA4 的 SFU 吞吐};
  \node[card, draw=ksTeal] (h3) at (68,-31)
    {高层架构设计\\
    数学改写 · Warp 分工 · 异步掩盖};
  \node[card, draw=ksTeal] (h4) at (0,-31)
    {编码实现\\
    按设计翻译为底层代码};

  \node[heading] (agent-title) at (34,-55)
    {当前编码 Agent：把算子当作 Code 工程任务};
  \node[card, draw=ksAmber] (a1) at (0,-68)
    {长程奖励稀疏\\
    变得短视，只求尽快跑通};
  \node[card, draw=ksAmber] (a2) at (68,-68)
    {跳过瓶颈分析与架构推演\\
    直接编写底层代码};
  \node[card, draw=ksAmber] (a3) at (68,-86)
    {缺少记录设计的结构化语言\\
    仅用 Markdown 零散记录};
  \node[card, draw=ksAmber] (a4) at (0,-86)
    {陷入恶性循环\\
    缺少设计 Taste，也缺少设计载体};

  \begin{scope}[on background layer]
    \node[frame, draw=ksTeal, fill=ksPaleTeal,
      fit=(human-title)(h1)(h2)(h3)(h4)] (human) {};
    \node[frame, draw=ksAmber, fill=ksPaleAmber,
      fit=(agent-title)(a1)(a2)(a3)(a4)] (agent) {};
  \end{scope}

  \draw[flow, ksTeal] (h1) -- (h2);
  \draw[flow, ksTeal] (h2) -- (h3);
  \draw[flow, ksTeal] (h3) -- (h4);
  \draw[flow, ksAmber] (a1) -- (a2);
  \draw[flow, ksAmber] (a2) -- (a3);
  \draw[flow, ksAmber] (a3) -- (a4);
  \draw[flow, ksInk, dashed] (human.south) --
    node[edge text, right=1mm] {根本差异：Design 先行，还是直接写 Code} (agent.north);
\end{tikzpicture}%
\endgroup%
%
\end{adjustbox}%
}{%
\MissingDiagram{04}%
}
\end{DiagramBlock}
现有编码 Agent 的做法正好相反。在当前的任务定义和强化学习训练方式下，Agent 得到的主要奖励是交出一份能编译、能通过测试的代码。长程任务中有效奖励稀疏，Agent 变得短视：跳过瓶颈分析，一上来就写底层代码。同时，Agent 也没有一种形式化的方式来记录、检查和修改设计，最多用 Markdown 写几句自然语言总结，没有结构约束，也无法检查前后是否自洽。结果是 Agent 一方面缺少对设计的敏感度和 Taste，另一方面缺少做设计的载体和能力。前者让它不去做设计，后者让它做了设计也留不下来，两者形成恶性循环。


\textbf{三个挑战的共同根源：人类把算子开发当作高层的软硬件协同 Design 问题，现有 Agent 则把它当作底层的 Code 任务。}设计一旦只能依附在代码上，三个环节的问题也因此出现：语料中的设计被工程实现细节淹没（挑战一），迭代中设计的好坏受实现质量牵连（挑战二），交付后设计取舍埋在代码里，难以审阅（挑战三）。


\subsection{总体思路：用 Kernel Schema 记录设计，接入迭代闭环}\label{sec:overall-approach}
既然根因在设计这一环，我们就从表达入手：在具体的 GPU 编程语言之上定义 Kernel Schema，只记录设计决策，不写语法细节和工程样板。\textbf{Kernel Schema 是对算子设计的压缩表示，用来准确记录、校验和维护设计，并接入自动化迭代 Loop。}


\begin{DiagramBlock}
\IfFileExists{diagrams/diagram-05.tex}{%
\begin{adjustbox}{max width=\linewidth,max totalheight=.80\textheight,center}
\begingroup%
\definecolor{ksBlue}{HTML}{3478B8}%
\definecolor{ksTeal}{HTML}{27877F}%
\definecolor{ksAmber}{HTML}{C58B28}%
\definecolor{ksInk}{HTML}{243447}%
\definecolor{ksSlate}{HTML}{667C92}%
\ifdefined\hypersetup\hypersetup{linkcolor=ksInk}\fi%
\begin{tikzpicture}[
  x=1mm, y=1mm,
  card/.style={draw, fill=white, rounded corners=2mm, line width=.7pt,
    text=ksInk, align=flush center, font=\fontsize{9.5}{12.5}\selectfont,
    inner xsep=2mm, inner ysep=1.7mm, outer sep=.35pt},
  chapter/.style={card, minimum width=112mm},
  connector/.style={draw=ksInk, line width=.7pt, rounded corners=2mm},
  flow/.style={connector, -{Stealth[length=2.25mm,width=1.65mm]}}
]
  \node[card, draw=ksSlate, minimum width=86mm] (root) at (0,0)
    {问题根源\\
    Agent 缺少设计敏感度，也缺少设计载体\\
    把 Design 问题当作 Code 任务处理};
  \node[card, draw=ksBlue, line width=.9pt, minimum width=86mm] (center) at (0,-21)
    {\textbf{方案：Kernel Schema}\\
    跨硬件、跨语言的算子设计压缩表示};

  \node[chapter, draw=ksTeal] (w1) at (6,-41)
    {第\ref{sec:schema}章 · 挑战一：数据与表示（已完成）\\
    算子压缩表示与知识蒸馏\\
    schema\_docs / template / extract / wiki};
  \node[chapter, draw=ksAmber] (w2) at (6,-61)
    {第\ref{sec:planner-coder}章 · 挑战二：生成与迭代（进行中）\\
    Planner–Coder 解耦闭环 · Step / Dispatch\\
    分离设计搜索与工程实现};
  \node[chapter, draw=ksSlate] (w3) at (6,-81)
    {第\ref{sec:graybox}章 · 挑战三：交付与维护（规划中）\\
    灰盒治理 · 分级异步验证 · 先验保序建模\\
    结合 Data 驱动与 Bench 验证};

  \coordinate (rail-top) at (-60,-21);
  \coordinate (rail-end) at (rail-top |- w3.west);
  \draw[flow] (root) -- (center);
  \draw[flow] (center.west) -- (rail-top) -- (rail-end) -- (w3.west);
  \draw[flow] (rail-top |- w1.west) -- (w1.west);
  \draw[flow] (rail-top |- w2.west) -- (w2.west);
\end{tikzpicture}%
\endgroup%
%
\end{adjustbox}%
}{%
\MissingDiagram{05}%
}
\end{DiagramBlock}
第\ref{sec:schema}章针对挑战一（数据与表示），介绍硬件感知、数据蒸馏与设计语言方面已完成的工作和实验结果；第\ref{sec:planner-coder}章针对挑战二（生成与迭代），介绍 Planner–Coder 分层优化闭环；第\ref{sec:graybox}章针对挑战三（交付与维护），介绍灰盒治理的规划与展望。


\SourceSeparator
\clearpage
\section{面向挑战一：Kernel Schema 压缩表示与知识蒸馏}\label{sec:schema}
本章针对数据与表示环节的问题：编程语言碎片化限制了 Agent 的设计思路，语料中的设计决策被工程代码掩盖，模型也缺少目标硬件的先验。我们为此定义了一种轻量的算子设计语言 \textbf{Kernel Schema}：它位于各种编程语言之上，只记录决定性能的设计决策。围绕它，我们做了四部分工作：语法规范（\InlineCode{schema\_\allowbreak{}docs}）、检查与压缩工具（\InlineCode{schema\_\allowbreak{}template}）、自动蒸馏管线（\InlineCode{schema\_\allowbreak{}extract}）和设计知识库（\InlineCode{schema\_\allowbreak{}wiki}）。


\begin{DiagramBlock}
\IfFileExists{diagrams/diagram-06.tex}{%
\begin{adjustbox}{max width=\linewidth,max totalheight=.80\textheight,center}
\begingroup%
\definecolor{ksBlue}{HTML}{3478B8}%
\definecolor{ksInk}{HTML}{243447}%
\begin{tikzpicture}[
  x=1mm, y=1mm,
  card/.style={draw=ksBlue, fill=white, rounded corners=2mm, line width=.7pt,
    text=ksInk, align=flush center, font=\fontsize{9.5}{12.5}\selectfont,
    minimum width=58mm, minimum height=20.5mm, inner xsep=2mm, inner ysep=1.7mm,
    outer sep=.35pt},
  flow/.style={-{Stealth[length=2.25mm,width=1.65mm]}, draw=ksBlue, line width=.7pt}
]
  \node[card] (d1) at (0,0) {1. 语言规范\\
    (schema\_docs)\\
    Contract / Design / Realization};
  \node[card] (d2) at (70,0) {2. 检查与压缩工具\\
    (schema\_template)\\
    合法性校验 · \textasciitilde{}10x 源码压缩};
  \node[card] (d3) at (70,-28) {3. 自动提取管线\\
    (schema\_extract)\\
    P1–P6 渐进蒸馏 · 盲审防幻觉};
  \node[card] (d4) at (0,-28) {4. 设计知识库\\
    (schema\_wiki)\\
    收录跨卡算子设计};

  \draw[flow] (d1) -- (d2);
  \draw[flow] (d2) -- (d3);
  \draw[flow] (d3) -- (d4);
\end{tikzpicture}%
\endgroup%
%
\end{adjustbox}%
}{%
\MissingDiagram{06}%
}
\end{DiagramBlock}
\SourceSeparator
\subsection{跨硬件、跨语言的算子设计语言（\InlineCode{schema\_\allowbreak{}docs}）}\label{sec:schema-docs}
我们希望这门语言在理论上能描述任意异构硬件（不同的硬件特性）和任意模型结构（不同的数学规格），并且描述完备。这首先是一个数据问题：只能从大量真实算子中归纳共性，再把共性压缩成少量字段，必须由数据驱动（Data-Driven）。CAKE-IR 同样以大量 Kernel 为数据归纳调度表达，但没有开源，且绑定专用的编译降级链。\textbf{Kernel Schema 比 CAKE-IR 这类编译器 IR 更高一层：不需要编译，只写决定性能的设计决策，Agent 有了描述和记录设计的语言，人类专家也能直接阅读。}


\InlineCode{schema\_\allowbreak{}docs} 的每个字段都推敲过原理和意义，没有一条是让 Agent 凭空撰写提出的。因为我们认为，至少目前，大模型还不擅长压缩，而压缩恰恰是智能的必要条件：它能写出冗长的规则，却很难判断哪些该删。因此我们采用人机结对的方式，以大量真实 SOTA 算子为样本，反复修订每一层抽象和每条字段规则：AI 从源码中拆出事实，专家判断这些事实的物理意义和取舍，再通过对话推演反例来修正规则。取舍上采用奥卡姆剃刀式的“删除测试”：一个字段或一句描述删掉后，如果既不会引起误读，也不会漏掉设计决策，就删去。


一份完整的 Kernel Schema 除了记录任务和来源的 \InlineCode{META\_\allowbreak{}DATA}，分为三个区域：只读的任务契约（\InlineCode{Contract}）、分三层描述的设计（\InlineCode{Design}），以及少量实现约束（\InlineCode{Realization}）。


\Needspace{6\baselineskip}
\subsubsection{\InlineCode{Contract}：任务要求与硬件规格（只读）}\label{sec:schema-contract}
一个算子任务，就是在给定的硬件上实现确定的数学输入输出。这两项在任务层面是固定的，\InlineCode{Contract} 区记录它们，对设计者和编码者都只读。


\textbf{软件算法要求}（\InlineCode{Algorithm\_\allowbreak{}requirement}）只从数学角度规定输入输出张量的规格、有效参数域、原位修改等副作用，以及以 PyTorch \InlineCode{reference + tolerance} 为准的误差判据，不包含任何实现假设。


\textbf{硬件平台规格}（\InlineCode{Hardware\_\allowbreak{}specification}）只登记会改变设计空间的执行层级和硬件特性，如 Cluster、TMEM、CTA-pair。更底层的硬件知识，我们不让 Agent 到算子源码库里去找：源码库只是未经蒸馏的实例，语法噪音多，会侵占上下文窗口，还会诱导 Agent 照搬旧卡的写法。这些知识被整理成分层 Wiki，包括架构能力（\InlineCode{Hardware}）、编程栈语义（\InlineCode{Programming}）、性能剖析（\InlineCode{Profile}）等。


分层 Wiki 中，\textbf{Microbench Wiki}（\InlineCode{Microbench}）提供定量数据。它在目标卡上对计算、访存和特殊函数单元做指令级实测，记录典型 Tile 规模下各硬件原语的实际周期与吞吐。有了这本定量账本，Agent 在设计阶段就能比较“一次矩阵乘 512 周期”和“一次 \InlineCode{exp2} 1024 周期”能否相互掩盖，不必等到编码阶段再去猜测或临时测量。


\Needspace{6\baselineskip}
\subsubsection{\InlineCode{Design}：数学变换、工作调度、微流水线三层决策}\label{sec:schema-design}
\InlineCode{Design} 区是 Schema 的主体，记录软硬件协同的设计决策，分为层层递进的三层，依次是数学变换（\InlineCode{Math\_\allowbreak{}reshape}）、工作调度（\InlineCode{Work\_\allowbreak{}schedule}）和微流水线（\InlineCode{Tile\_\allowbreak{}pipeline}）。


\begin{DiagramBlock}
\IfFileExists{diagrams/diagram-07.tex}{%
\begin{adjustbox}{max width=\linewidth,max totalheight=.80\textheight,center}
\begingroup%
\definecolor{ksBlue}{HTML}{3478B8}%
\definecolor{ksTeal}{HTML}{27877F}%
\definecolor{ksAmber}{HTML}{C58B28}%
\definecolor{ksInk}{HTML}{243447}%
\begin{tikzpicture}[
  x=1mm, y=1mm,
  card/.style={draw, fill=white, rounded corners=2mm, line width=.8pt,
    text=ksInk, align=flush center, font=\fontsize{9.5}{12.5}\selectfont,
    minimum height=20mm, inner xsep=2mm, inner ysep=1.7mm, outer sep=.4pt},
  flow/.style={-{Stealth[length=2.6mm,width=1.9mm]}, draw=ksInk, line width=.9pt},
  edge text/.style={text=ksInk, font=\fontsize{9}{12.5}\selectfont, inner sep=1mm}
]
  \node[card, draw=ksBlue, minimum width=100mm] (l1) at (0,0)
    {\textbf{第一层：Math\_reshape（数学变换）}\\
    等价改写计算顺序 · 在线状态递推 · 精度取舍\\
    难度最高，可省下的计算量与访存最多};
  \node[card, draw=ksTeal, minimum width=120mm] (l2) at (0,-33)
    {\textbf{第二层：Work\_schedule（Tile 以上的全局调度）}\\
    同构 Worktile 切分、SM / CTA / Cluster 映射与负载均衡\\
    难度中等，决定设备级并行利用率};
  \node[card, draw=ksAmber, minimum width=140mm] (l3) at (0,-66)
    {\textbf{第三层：Tile\_pipeline（Tile 内微流水线）}\\
    Microbench 周期先验 · Warp 角色划分 · 8 动词与等待边编排异步重叠\\
    问题最具体，决定能多接近硬件峰值吞吐};

  \draw[flow] (l1) -- node[edge text, right=1mm]
    {确定计算哪些量、跨轮携带哪些状态} (l2);
  \draw[flow] (l2) -- node[edge text, right=1mm]
    {确定单个 Worker 的子任务边界} (l3);
\end{tikzpicture}%
\endgroup%
%
\end{adjustbox}%
}{%
\MissingDiagram{07}%
}
\end{DiagramBlock}
三层构成一座决策金字塔：从上到下，抽象粒度逐层变细，\textbf{优化难度和可能带来的性能收益逐层降低}。三层都不写代码，不涉及 Host 启动、IO 声明、内存分配和底层 Layout 代数，记录的只是抽象层面的设计决策。


\textbf{第一层} \InlineCode{Math\_\allowbreak{}reshape}\textbf{（数学变换）}在纯数学层面做等价重排或有界的精度折中，如分块矩阵乘、Online Softmax、线性注意力的逐块三角逆展开，目的是减少计算量、缩短串行依赖，或省去全局内存中的中间结果；收益与代价记在 \InlineCode{Ledger} 账本中核对。

例如 \InlineCode{FA4 Prefill} 用滞后阈值 $\tau=4$ 跳过不必要的重缩放，分子窄化为 FP8，分母保留未舍入的 FP32 累加，十余行就能写清楚：
\begin{CodeBox}
\begin{Verbatim}[breaklines,breakanywhere,breaksymbolleft={},breaksymbolright={},breakautoindent=true]
# 摘自 schema_template/examples/fa4_prefill/root.yaml (DESIGN.Math_reshape)
- Name: online
  Math: |
    keep=(c*(mhat(i-1)-safe)>=-tau); mhat(i)=mhat(i-1) if keep, otherwise a; base=mhat(i-1) if keep, otherwise safe.
    alpha(i)=1 if keep, otherwise exp2_fast(c*(mhat(i-1)-safe)); stats(i)[r,:]=(mhat(i-1),base).
    exponentials(i)[r,k]=exp2_fast(fma32(c,maskedscore(i)[r,k],offset-c*base)); prob(i)=fp8_e4m3fn(exponentials(i)).
    ell(i)=alpha(i)*ell(i-1)+sum32_k exponentials(i)[r,k]; numerator(i)=mma32(prob(i),vc(i),alpha(i)*numerator(i-1)).
    invell=1/ell(tile_count-1) if ell is nonzero and not NaN, otherwise zero; oblock=bf16(numerator(tile_count-1)*invell).
  Ledger:
    Baseline: FP32 softmax weights in both numerator and denominator
    Gain: FP8 PV with bounded per-work online state
    Cost: FP8 numerator rounding, lagged maximum (tau=4) and approximate exponentials; denominator uses unrounded exponentials
\end{Verbatim}
\end{CodeBox}

\textbf{第二层} \InlineCode{Work\_\allowbreak{}schedule}\textbf{（工作调度）}负责 Tile 以上的并行规划：把改写后的任务切分为可统一处理的子任务单元（\InlineCode{Worktile}，其中 \InlineCode{Invariant} 在块内常驻，\InlineCode{Iteration} 逐轮流入），确定归并方式，映射到硬件资源（如单 CTA 或 2-CTA Cluster），再确定任务顺序和负载均衡方式（静态步长常驻，或原子票据动态队列）。例如，\InlineCode{FA4 Prefill} 的工作调度如下：
\begin{CodeBox}
\begin{Verbatim}[breaklines,breakanywhere,breaksymbolleft={},breaksymbolright={},breakautoindent=true]
# 摘自 schema_template/examples/fa4_prefill/packed_persistent.yaml (DESIGN.Work_schedule)
Decomposition:
  Partition: {cluster_tile: {size: cpad, tile: 1, meaning: cluster tile coordinate}}
  Worktile:
    Invariant: {qc: {shape: [{row: packedrows}, {channel: channels}], dtype: fp8_e4m3fn, meaning: qc operand block}}
    Iteration:
      kc: {shape: [{key: kvblock}, {channel: channels}], dtype: fp8_e4m3fn, direction: ascending, meaning: kc operand block}
      vc: {shape: [{key: kvblock}, {channel: channels}], dtype: fp8_e4m3fn, direction: ascending, meaning: vc operand block}
Mapping:
  Worker: {unit: cta, unit_count: 2}          # 2-CTA 协作构成单个逻辑 Worker
  Placement: {unit: cluster, unit_count: 1}
  Worker_count: min(74,cpad)                  # 常驻 74 组 CTA Pair 复用流水状态
Schedule:
  Worker_assignment: {Preassigned: {start: worker_id, stride: worker_count}, Refill: none, Exit: assigned_done}
\end{Verbatim}
\end{CodeBox}

\textbf{第三层} \InlineCode{Tile\_\allowbreak{}pipeline}\textbf{（微流水线）}描述单个 Tile 内部的执行。它参考 \InlineCode{Microbench} 的周期数据，把 Warp / Warpgroup 划分为专职角色，用 8 个跨平台的基础动词和 3 类等待契约描述片上存储复用与指令级并发掩盖。8 个动词按用途分为数值计算 \InlineCode{compute /\allowbreak{} matmul}、搬运访存 \InlineCode{load /\allowbreak{} store}、顺序同步 \InlineCode{await /\allowbreak{} signal}、资源占用 \InlineCode{acquire /\allowbreak{} release}；3 类等待契约是数据就绪 \InlineCode{Prepared}、槽位释放 \InlineCode{Released} 和协议放行 \InlineCode{Allowed}。

下面节选了 \InlineCode{FA4 Prefill} 中 MMA 角色的片段：
\begin{CodeBox}
\begin{Verbatim}[breaklines,breakanywhere,breaksymbolleft={},breaksymbolright={},breakautoindent=true]
# 摘自 schema_template/examples/fa4_prefill/packed_persistent.yaml (DESIGN.Tile_pipeline · MMA 角色)
Pipe:
  - for(t):
      - load:work(t)
      - if(valid(t) and tile_count(t) > 0):
          - acquire:scoreslot(n[t,0])
          - for(d): [matmul:qk(n[t,0],d)]
          - for(n[t,1:]):                     # 稳态主循环：当前块 QK 与上一块 PV 交错推进
              - acquire:scoreslot(n)
              - for(d): [matmul:qk(n,d)]
              - for(d): [matmul:pv(n-1,d)]
          - for(d): [matmul:pv(n[t,tile_count(t)-1],d)]
Waiting:
  Prepared:
    - '{issue matmul:qk(n,d)} ← [{done LOAD_QK[*]::load:qc(t,d)}, {done LOAD_QK[*]::load:kc(n,d)}]'
    - '{issue matmul:pv(n,d)} ← [{done SOFTMAX[*]::compute:prob(n)}, {done LOAD_V[*]::load:vc(n,d)}]'
  Released:
    - '{done acquire:scoreslot(n+2)} ← [{done SOFTMAX[*]::load:score(n)}]'   # 双槽缓冲安全回绕条件
\end{Verbatim}
\end{CodeBox}
\Needspace{6\baselineskip}
\subsubsection{\InlineCode{Realization}：规定必须使用或禁止的原语}\label{sec:schema-realization}
\InlineCode{Realization} 区只作少量补充：用 \InlineCode{Require /\allowbreak{} Prefer /\allowbreak{} Forbid} 规定少数决定方案能否成立的原语或预算边界，例如在 Blackwell 上要求使用 \InlineCode{tcgen05} 和 TMA、禁止退回 \InlineCode{mma.\allowbreak{}sync}，不涉及一般的代码组织。


\newpage

\subsection{形式化检查、篇幅压缩与流水线验证（\InlineCode{schema\_\allowbreak{}template}）}\label{sec:schema-template}
语言定义之外，还需要工具来检查 Agent 写出的 Schema：结构是否合法，格式是否统一，篇幅是否精炼，流水线有没有死锁。\InlineCode{schema\_\allowbreak{}template} 提供检查、规范化和可视化工具（入口为 \InlineCode{check\_\allowbreak{}yaml.\allowbreak{}py}），并附有 \InlineCode{FA4 Prefill}、\InlineCode{FA4 Decode}、\InlineCode{GDN Prefill}、\InlineCode{QSA Prefill} 四个算子的标准样例。主要功能和结果有以下三项。


\begin{enumerate}[start=1]
\item
\textbf{结构校验与篇幅上限}：自动校验字段结构、跨区符号和属性绑定，并幂等地统一大小写、引号与格式。每类字段都有字数和行数上限，如三层 \InlineCode{Thesis} 不超过 2 行、单条 \InlineCode{Math} 不超过 10 行、动作说明不超过 1 行，防止 Agent 堆砌文字。单份完整设计因此控制在 300～500 行。


\item
\textbf{压缩比}：\InlineCode{FA4 Prefill} 的主核代码有 4,040 行（加上归并与接口模块超过 8,800 行），对应的 \InlineCode{root.\allowbreak{}yaml} 只有 384 行，压缩比 10.5×；\InlineCode{GDN Prefill} 的可读源码 5,938 行，完整 Schema 482 行，压缩比 12.3×。几百行的 Schema 就是 Agent 的设计记忆：数学改写、工作调度和流水编排这些决定性能的决策都在里面，不必再到数千行代码里去找。


\item
\textbf{流水线验错与依赖图}：\InlineCode{Tile\_\allowbreak{}pipeline} 中的动作和等待有形式化语义。\InlineCode{check\_\allowbreak{}yaml.\allowbreak{}py} 在给定的样本参数下展开代表性 Worker 的循环与分支，构建 \InlineCode{issue /\allowbreak{} done} 事件有向图，自动检查重复动作、悬空引用和同步死锁环，并渲染成流水线依赖图（见下图）。


\end{enumerate}

\FullPageImage{assets/root_pipeline.png}

\newpage

\subsection{自动蒸馏：从 \InlineCode{kernel\_\allowbreak{}wiki} 到 \InlineCode{schema\_\allowbreak{}wiki}（\InlineCode{schema\_\allowbreak{}extract}）}\label{sec:schema-extract}
有了语言和检查工具，下一步是“造数据”：从原始算子库中批量清洗、提取专家的设计，整理成设计知识库 \InlineCode{schema\_\allowbreak{}wiki}，供 Agent 在设计新算子或适配新卡时检索。


\OriginalImage{assets/kernel_wiki.png}{Kernel Wiki 算子题库与验收看板}{\emph{图：标准算子评测与源码题库（Kernel Wiki）看板。覆盖 Dense Attention（16）、Sparse Attention（8）、Linear Attention（15）、GEMM（17）、MoE（16）五类共 72 道动态支持域题目（含 720 个固定评测 Shape 与统一的容差与计时规范），用作蒸馏与闭环验证的题库。}}
基于这 72 道题，我们把 L20D（B300）主批次中验证证据完整的 44 个算子（20 个 AKA 优胜解 + 24 个参考实现）导入 \InlineCode{kernel\_\allowbreak{}wiki}，并实现了离线自动蒸馏工具 \textbf{Schema Extract}（\InlineCode{schema\_\allowbreak{}extract}，入口为 \InlineCode{extract\_\allowbreak{}yaml.\allowbreak{}py}），流程如下图。


\begin{DiagramBlock}
\IfFileExists{diagrams/diagram-08.tex}{%
\begin{adjustbox}{max width=\linewidth,max totalheight=.80\textheight,center}
\begingroup%
\def\baselinestretch{1}%
\definecolor{D08Blue}{HTML}{3478B8}%
\definecolor{D08Teal}{HTML}{27877F}%
\definecolor{D08Amber}{HTML}{C58B28}%
\definecolor{D08Ink}{HTML}{243447}%
\begin{tikzpicture}[
  x=1mm, y=1mm,
  card/.style={draw, fill=white, rounded corners=2mm, line width=.75pt,
    minimum width=75mm, minimum height=13mm, inner xsep=3mm, inner ysep=2.5mm,
    outer sep=.4pt, align=flush center, text=D08Ink, font=\fontsize{9.5}{12.5}\selectfont},
  flow/.style={draw=D08Ink, line width=.7pt, rounded corners=1.5mm,
    -{Latex[length=2.5mm,width=1.6mm]}},
  guard flow/.style={flow, draw=D08Amber, dash pattern=on 2.6pt off 2pt}
]
\node[card, draw=D08Blue] (src) at (-41,0) {原始算子与 Contract\\
  （kernel\_wiki · L20D 44 题）};
\node[card, draw=D08Teal] (p12) at (-41,-23) {P1 建立共享 Contract\\
  P2 追踪调用链，识别代表机制};
\node[card, draw=D08Teal] (p34) at (-41,-46) {P3 提取 Root 三层设计\\
  P4 规划多路径方案};
\node[card, draw=D08Teal] (p56) at (-41,-69) {P5 构建附加设计\\
  P6 全算子盲审与原子发布};
\node[card, draw=D08Blue] (out) at (-41,-92) {设计知识库\\
  （schema\_wiki · 已入库 18 题 / 48 份设计）};

\node[card, draw=D08Amber, minimum width=68mm] (guard) at (44,-46)
  {源码行号与引文哈希核验\\
   check\_yaml 规范检查\\
   独立 Reviewer 会话盲审};

\draw[flow] (src) -- (p12);
\draw[flow] (p12) -- (p34);
\draw[flow] (p34) -- (p56);
\draw[flow] (p56) -- (out);
\coordinate (guard-upper) at (0,-41);
\coordinate (guard-lower) at (0,-51);
\draw[guard flow] (guard.west |- guard-upper) -- (guard-upper)
  -- (guard-upper |- p12.east) -- (p12.east);
\draw[guard flow] (guard.west) -- (p34.east);
\draw[guard flow] (guard.west |- guard-lower) -- (guard-lower)
  -- (guard-lower |- p56.east) -- (p56.east);
\end{tikzpicture}%
\endgroup%
%
\end{adjustbox}%
}{%
\MissingDiagram{08}%
}
\end{DiagramBlock}
\textbf{分阶段提取与防幻觉校验。}提取过程拆成 P1～P6 六个阶段，每个阶段在独立会话中完成。每个阶段，模型都要提交结构化证据，标明对应的源码行号和原文片段，由 Runner 核对文件哈希与原文。产出经 \InlineCode{check\_\allowbreak{}yaml.\allowbreak{}py} 规范化后，交给一个新开的只读 Reviewer 会话盲审，检查脱离源码能否读懂、分层是否准确、有无冗余或遗漏，未通过就打回。P6 终审通过后原子发布，并生成加密回执（\InlineCode{receipt.\allowbreak{}json}）。


\textbf{所用模型与当前规模。}在上述约束下，直接使用国产模型 \InlineCode{qwen3.\allowbreak{}8-\allowbreak{}max} 就能自动完成复杂算子的提取并通过盲审。目前已在 L20D 44 题主库上完成首批 18 个复杂算子、共 48 份设计的入库。例如 \InlineCode{bf16\_\allowbreak{}nvfp4\_\allowbreak{}gemm} 的多路径候选源码共 18,450 行，提取后得到 7 份正交设计，共 1,887 行。


\textbf{检索到的是设计思路。}\InlineCode{schema\_\allowbreak{}wiki} 回答的是两类问题：同一张卡上的其他算子是怎么设计的，同一个算子在其他卡上又是怎么设计的。它不教 CuTe 等 DSL 的写法，那是 Agent 本来就擅长的部分。

\newpage

\subsection{初步验证：无外部参考的 Schema 注入消融实验}\label{sec:schema-ablation}
这组实验检验一个问题：只在提示词中提供 Kernel Schema，能否改善编码 Agent 的生成与收敛。实验不使用第\ref{sec:planner-coder}章的双层架构，只做最简单的注入（Prompt-only Schema Injection）。


\textbf{实验设置。}选取 GDN Prefill（17 shapes）、FA4 Prefill（30 shapes）、FA4 Decode（30 shapes）、QSA Prefill（4 shapes）四个算子，在 Agate L20D（SM103）上用 \InlineCode{CuTeDSL} 从零编写。实验关闭所有外部 Wiki、开源仓库和联网检索，在相同的简化版 AKA-light 循环下运行 4 算子 × 2 组（Schema vs. Control）× 3 次重复，共 24 路隔离对照。两组的唯一区别是：\InlineCode{Schema} 组每轮提示词中附带一份简短的 Kernel Schema，要求 Agent 按这份设计实现，相当于事先拿到了解题思路；\InlineCode{Control} 组没有这份设计，自行探索。

\OriginalImage{assets/01_early_convergence.png}{早期收敛：同等 Token 谁更快，同等性能谁更省}{\emph{图：前 4M 累计输出 Token 内的早期收敛轨迹与达标预算（蓝色为} \InlineCode{Schema} \emph{组，橙色为} \InlineCode{Control} \emph{组）。纵轴为对数延迟，越低越快。}}
\OriginalImage{assets/03_best_performance.png}{收敛效率与最佳性能汇总}{\emph{图：A 行为固定性能目标下的首达 Token 消耗；B 行为固定 4M / 6M / 8M 输出 Token 预算下的最佳延迟与中位加速比；C 行为共同迭代窗口末的分布与极值。}}
实验结果归纳为以下两点。


\begin{enumerate}[start=1]
\item
\textbf{早期收敛更快：达到同一门槛，Token 少用约 40\%～50\%}


有了设计蓝图，Agent 在迭代早期走的弯路更少：


\begin{itemize}
\item
\textbf{GDN Prefill}：以 ≤200 µs 为门槛，\InlineCode{Schema} 组 3/3 达标，最快两路只用了 1.75M / 2.17M 输出 Token；\InlineCode{Control} 组 2/3 达标，首次达标用了 3.68M / 3.76M，Schema 组约省一半预算。


\item
\textbf{QSA Prefill}：前 1M 输出 Token 内，\InlineCode{Schema} 组已有 2/3 产出被接受的有效解，\InlineCode{Control} 组为 0/3；门槛收紧到 ≤350 µs 时，达标率为 3/3 对 1/3。


\item
\textbf{FA4 Decode}：以 ≤40 µs 为门槛，达标率为 3/3 对 2/3；前 2M Token 内，\InlineCode{Schema} 组的延迟曲线下降更陡。


\end{itemize}
\item
\textbf{各 Run 结果更集中，中位延迟更低}


衡量指标是几何标准差 GSD（越接近 1.0，各 Run 结果越一致）和同等 Token 预算下的中位加速比。三次重复都按同一份设计推进，不容易各自走散。


\begin{itemize}
\item
\textbf{GDN Prefill}：中位、最慢、最快结果都好于 \InlineCode{Control} 组。4M / 6M / 8M 预算下中位加速比均为 1.21×～1.22×；共同窗口末最快解加速 1.36×（122.5 µs vs 166.3 µs，之后进一步降至 108.4 µs）；GSD 从 1.57 降至 1.17，\InlineCode{Control} 组中长期卡在 365.9 µs 的那条掉队慢路，在 \InlineCode{Schema} 组没有出现。


\item
\textbf{FA4 Decode}：6M～8M 预算下中位加速 1.09×～1.11×；后段 GSD 从 1.17 降至 1.04，\InlineCode{Control} 组中停在 46.6 µs 的慢路在 \InlineCode{Schema} 组没有出现。


\item
\textbf{QSA Prefill}：4M～8M 预算下中位加速 1.16×～1.30×，GSD 从 2.50 降至 1.01。


\end{itemize}
\end{enumerate}

\begin{ReportQuote}
\textbf{小结}：\textbf{只在提示词中加入一份数百行的 Kernel Schema，Agent 就能更早达到性能门槛，各 Run 结果也更集中。}GDN Prefill 达标少用约一半 Token；GDN Prefill、FA4 Decode、QSA Prefill 的 GSD 分别降至 1.17、1.04、1.01。

这里的 Schema 是固定的，Agent 只负责照着实现。下一章的 Planner–Coder 分层闭环更进一步：Planner 搜索和改进 Schema，Coder 忠实实现并反馈。


\end{ReportQuote}
\SourceSeparator
\clearpage


\section{面向挑战二：Planner–Coder 两阶段闭环}\label{sec:planner-coder}
第\ref{sec:schema}章讨论了如何压缩表示算子设计，本章讨论如何在迭代中使用 Schema，应对第~\ref{sec:challenge-iteration}~节的几个现象：多次运行的结果离散，设计与实现相互牵连，试错会消耗大量时间和 Token，稀疏奖励又让 Agent 只顾眼前。这些现象指向同一个做法：先解耦，再分阶段处理。我们的做法是\textbf{把策略思考和代码修改放进两个阶段，由 Planner 和 Coder 分别负责}。


\SourceSeparator
\subsection{思路：用两个角色分开设计与实现}\label{sec:two-roles}
现有报道的一些工作，倾向于通过加重循环编排来提升效果。例如 Kernel Forge 用 MCTS 树搜索探索多条优化路径，并配套模型算子捕获和带回退的部署封装；另一些系统设置总控、测试、规划、编码、合并等多个角色，甚至嵌套多级 RSI。我们可以把这类做法统称为 \textbf{Loop Engineering}：效果主要靠工程编排的堆叠，系统随之变得非常复杂。它可以扩大搜索覆盖面，但层级和角色增加后，角色之间传递信息会有损耗，状态管理的开销上升，每个角色上下文里与自身任务相关的内容比例也会下降。


相比之下，\textbf{Humanize} 这种最小化的轻量系统更可取：只设两个职责正交的角色，一个独立审视当前结果并提出高层改进意见，另一个负责把意见落地执行。两个角色就能组成闭环，不需要复杂的编排。而且，它也取得了比较好的效果；因此我们也认为，理想的系统设计应该遵从轻量、最小化、Less is more 的设计原则。


\begin{DiagramBlock}
\IfFileExists{diagrams/diagram-09.tex}{%
\begin{adjustbox}{max width=\linewidth,max totalheight=.80\textheight,center}
\begingroup%
\def\baselinestretch{1}%
\definecolor{D09Blue}{HTML}{3478B8}%
\definecolor{D09Teal}{HTML}{27877F}%
\definecolor{D09Amber}{HTML}{C58B28}%
\definecolor{D09Ink}{HTML}{243447}%
\definecolor{D09PaleBlue}{HTML}{F0F6FC}%
\definecolor{D09PaleAmber}{HTML}{FFF8E8}%
\begin{tikzpicture}[
  x=1mm, y=1mm,
  card/.style={draw, fill=white, rounded corners=2mm, line width=.75pt,
    inner xsep=2.5mm, inner ysep=2.6mm, outer sep=.4pt, align=flush center,
    text=D09Ink, font=\fontsize{9.5}{12.5}\selectfont},
  group box/.style={draw, rounded corners=2.5mm, line width=.7pt,
    inner sep=0pt, outer sep=.4pt},
  group title/.style={inner sep=0pt, text=D09Ink, align=flush center,
    font=\bfseries\fontsize{10}{12.5}\selectfont},
  edge text/.style={fill=white, inner xsep=1mm, inner ysep=.5mm, align=flush center,
    text=D09Ink, font=\fontsize{9}{12.5}\selectfont},
  flow/.style={draw=D09Ink, line width=.7pt, rounded corners=1.8mm,
    -{Latex[length=2.5mm,width=1.6mm]}},
  schema flow/.style={flow, draw=D09Blue, line width=.9pt,
    -{Latex[length=3mm,width=2mm]}}
]
\node[group title] at (0,17) {原生 AKA：单层混合循环，缺少高层设计约束};
\node[group title] at (0,-48)
  {两阶段解耦：Planner 负责压缩态设计，Coder 负责解压与实现};

\node[card, draw=D09Amber, minimum width=42mm, minimum height=18mm] (o1) at (-55,0)
  {从零开始 / 基于已有 SOTA};
\node[card, draw=D09Amber, minimum width=46mm, minimum height=18mm] (o2) at (0,0)
  {单一 Episode 内混合迭代\\
   策略、代码与编译排错交织};
\node[card, draw=D09Amber, minimum width=46mm, minimum height=18mm] (o3) at (55,0)
  {验证与晋升判定\\
   易因实现受阻\\
   放弃正确方向};

\node[card, draw=D09Blue, minimum width=51mm, minimum height=24mm] (n1) at (-48,-71)
  {Planner 维护 Schema\\
   Design / Extract 生成初版\\
   Optimize 分层演进};
\node[card, draw=D09Teal, minimum width=55mm, minimum height=24mm] (n2) at (48,-71)
  {Coder 维护 Kernel（基于 AKA）\\
   Implement 将 Schema 解压为代码\\
   内部处理编译与工程细节};

\begin{scope}[on background layer]
  \node[group box, draw=D09Amber, fill=D09PaleAmber,
    fit={(-81,23) (81,-25)}] (old) {};
  \node[group box, draw=D09Blue, fill=D09PaleBlue,
    fit={(-81,-41) (81,-88)}] (new) {};
\end{scope}

\draw[flow] (o1) -- (o2);
\draw[flow] (o2) -- (o3);
\coordinate (retry) at (0,-18);
\draw[flow, draw=D09Amber] (o3.south) -- (o3.south |- retry)
  -- node[edge text, fill=D09PaleAmber] {缺少方向的重试} (retry) -- (o2.south);

\coordinate (schema) at (0,-65);
\coordinate (evidence) at (0,-77);
\draw[schema flow] (n1.east |- schema) --
  node[edge text, fill=D09PaleBlue] {下发 Schema\\限定方向与边界} (n2.west |- schema);
\draw[schema flow, draw=D09Teal] (n2.west |- evidence) --
  node[edge text, fill=D09PaleBlue] {回传实测结果\\不含编译排错细节} (n1.east |- evidence);

\draw[flow, draw=D09Blue, dash pattern=on 2.6pt off 2pt] (old.south) --
  node[edge text, right=2mm] {引入 Kernel Schema，拆分职责} (new.north);
\end{tikzpicture}%
\endgroup%
%
\end{adjustbox}%
}{%
\MissingDiagram{09}%
}
\end{DiagramBlock}
现有的 \textbf{AKA（Atrex-Kernel-Agent）} 从零开始或从已有 SOTA 算子出发，在每个 Episode 内反复修改代码、编译评测，并晋升更好的版本。作为工具，AKA 已相当成熟：它完整实现了自闭环，Agent 可以在里面持续尝试和改进算子。问题在于缺少高层设计约束，它实际上是一个无约束、无方向的 Optimize 与 Implement 循环，设计和实现混在同一个 Episode 里。


因此，我们复用 AKA 的底层执行闭环，以 \textbf{Kernel Schema} 作为两者之间的交接物，拆出两个角色：


\begin{itemize}
\item
\textbf{Planner（设计者，操作压缩态的 Schema）}：不接触底层语法和编译细节，在百行量级的设计描述上决定优化方向；


\item
\textbf{Coder（实现者，操作解压态的 Kernel）}：由 AKA 改造而来，在 Schema 规定的范围内把设计还原为高性能代码。


\end{itemize}
\SourceSeparator
\subsection{视角：为什么需要在设计空间里搜索}\label{sec:design-space-search}
第~\ref{sec:challenge-iteration}~节把迭代难以收敛归因于设计与实现耦合。换一个角度，这也可以从两个搜索空间的性质来理解。直接迭代代码，相当于在代码空间里做局部搜索：每一步改动若干行，编译评测后决定是否保留。这个空间维度很高，而且很不连续，一处索引写错、一个屏障漏配，结果就从“快”跳到“错”或“超时”。

\textbf{借用生物遗传学中的适应性景观（Fitness Landscape）的说法，算子自身的Kernel空间中路径崎岖弯折、存在许多不收敛的病态点，仅一次评测或NCU反馈信息很难告诉 Agent 下一步该往哪里走}。但是，Schema 描述的设计空间要小得多：一份设计由计算路径、分块与映射、流水角色和缓冲槽数等百行量级的决策构成。直观地说，相邻设计之间的性能变化通常有迹可循，\InlineCode{Microbench} 的信息还能在写代码之前给出粗略的估计。
\OriginalImage{assets/fitness_landscape.png}{两种搜索空间的适应性景观示意}{\emph{图：左为 Kernel 代码空间，景观崎岖、断崖密布，从同一起点出发的几条搜索路径各自停在不同的局部峰上；右为 Schema 设计空间，景观平滑，路径顺坡即可攀升到主峰或次峰。示意图，不对应具体数据。}}

两个空间之间可以理解为一种一对多的映射：一份设计对应许多份实现，Schema 是这些实现的压缩，代码是 Schema 的展开。问题在于，这种映射并不保持邻域关系。设计空间里的一小步，例如把 \InlineCode{FA4 Prefill} 的 score 缓冲由双槽改为三槽，或让当前块的 QK 与上一块的 PV 交错推进，在 Schema 中只是改动 \InlineCode{Tile\_\allowbreak{}pipeline} 的几行循环和等待条件；落到代码里，却要重新安排存储分配、屏障相位和软件流水的首尾，常常牵动成百上千行。反过来，代码里只差一处的两份实现，可能一份正确且接近设计上限，另一份死锁或读到未就绪的数据；挪动一处等待的位置，也可能让两份文本相近的代码实际执行着不同的流水设计。

\textbf{性能主要由设计决定，而代码空间里的“邻居”与性能上的远近并不对应，局部搜索因此很难找准优化的方向。}


\begin{DiagramBlock}
\IfFileExists{diagrams/diagram-14.tex}{%
\begin{adjustbox}{max width=\linewidth,max totalheight=.80\textheight,center}
\begingroup%
\def\baselinestretch{1}%
\definecolor{D14Blue}{HTML}{3478B8}%
\definecolor{D14Teal}{HTML}{27877F}%
\definecolor{D14Amber}{HTML}{C58B28}%
\definecolor{D14Ink}{HTML}{243447}%
\definecolor{D14PaleBlue}{HTML}{F0F6FC}%
\definecolor{D14PaleAmber}{HTML}{FFF8E8}%
\begin{tikzpicture}[
  x=1mm, y=1mm,
  card/.style={draw, fill=white, rounded corners=2mm, line width=.75pt,
    inner xsep=2.5mm, inner ysep=2.6mm, outer sep=.4pt, align=flush center,
    text=D14Ink, font=\fontsize{9.5}{12.5}\selectfont},
  group box/.style={draw, rounded corners=2.5mm, line width=.7pt,
    inner sep=0pt, outer sep=.4pt},
  group title/.style={inner sep=0pt, text=D14Ink, align=flush center,
    font=\bfseries\fontsize{10}{12.5}\selectfont},
  edge text/.style={fill=white, inner xsep=1mm, inner ysep=.5mm, align=flush center,
    text=D14Ink, font=\fontsize{9}{12.5}\selectfont},
  flow/.style={draw=D14Ink, line width=.7pt, rounded corners=1.8mm,
    -{Latex[length=2.5mm,width=1.6mm]}},
  expand/.style={flow, draw=D14Teal, dash pattern=on 2.6pt off 2pt}
]
\node[group title] at (0,17)
  {设计空间（Schema，压缩态）：维度低，邻域较平滑};
\node[group title] at (0,-31)
  {代码空间（Kernel，解压态）：维度高，邻域崎岖};

\node[card, draw=D14Blue, minimum width=46mm, minimum height=16mm] (da) at (-52,0)
  {设计 A\\双槽缓冲，QK 与 PV 串行};
\node[card, draw=D14Blue, minimum width=46mm, minimum height=16mm] (db) at (52,0)
  {设计 B\\三槽缓冲，QK 与 PV 交错};

\node[card, draw=D14Teal, minimum width=46mm, minimum height=16mm] (a1) at (-52,-50)
  {实现 A-1\\正确，接近设计 A 的上限};
\node[card, draw=D14Amber, minimum width=46mm, minimum height=16mm] (a2) at (-52,-78)
  {实现 A-2\\与 A-1 只差一处屏障\\结果错误或超时};
\node[card, draw=D14Teal, minimum width=46mm, minimum height=16mm] (b1) at (52,-64)
  {实现 B-1, B-2, …\\重排角色、屏障与存储分配};

\begin{scope}[on background layer]
  \node[group box, draw=D14Blue, fill=D14PaleBlue,
    fit={(-81,23) (81,-13)}] (design) {};
  \node[group box, draw=D14Amber, fill=D14PaleAmber,
    fit={(-81,-25) (81,-92)}] (code) {};
\end{scope}

\draw[flow, draw=D14Blue, line width=.9pt] (da) --
  node[edge text, fill=D14PaleBlue] {一小步\\改动数行 Schema} (db);

\draw[expand] (da.south) -- (a1.north);
\draw[expand] (db.south) -- (b1.north);
\node[edge text] at (-52,-19) {一对多展开};
\node[edge text] at (52,-19) {一对多展开};

\draw[flow, draw=D14Amber] (a1.east) -- (-8,-50) -- (-8,-64) -- (b1.west);
\node[edge text, fill=D14PaleAmber] at (10.5,-57) {一大步\\成百上千行改动};
\draw[draw=D14Amber, line width=.7pt, dash pattern=on 2.6pt off 2pt] (a1.south) --
  node[edge text, fill=D14PaleAmber, right=2mm] {代码相近，性能悬殊} (a2.north);
\end{tikzpicture}%
\endgroup%
%
\end{adjustbox}%
}{%
\MissingDiagram{14}%
}
\end{DiagramBlock}
这也解释了 Planner 与 Coder 的分工。Planner 在维度低、较平滑的设计空间里走“大步”，比较、选择和修改设计；Coder 只在一份设计所对应的代码邻域里做局部实现与调优，也就是第~\ref{sec:design-tree-workflow}~节 Implement 环节所说的解压与信息恢复。设计每迈出一步，都由 Coder 在新的代码邻域里把实现补全、调稳，评测结果因此更多反映设计本身的上限，较少受某一次实现瑕疵的干扰。同一节介绍的设计遗传树，记录的正是设计空间里的这些“Stride”。


\SourceSeparator
\subsection{工作流程：两种起点、分层优化与设计遗传树}\label{sec:design-tree-workflow}
两阶段架构把算子开发拆成 \textbf{Design}、\textbf{Extract}、\textbf{Implement}、\textbf{Optimize} 四个动作，并用同一套流程处理“从零构建”和“从已有 SOTA 继续优化”两种情况。


初始 Schema 有两种来源。从零开始（From Scratch）时，Planner 执行 \textbf{Design}，参考任务的 \InlineCode{Contract} 和 \InlineCode{schema\_\allowbreak{}wiki} 中相近算子的设计，写出第一版 Schema；从已有 SOTA Kernel 开始时，Planner 执行 \textbf{Extract}，从 SOTA 源码中提取出与之一致的初始 Schema。


有了初始 Schema 之后，Implement 与 Optimize 交替进行：


\begin{itemize}
\item
\textbf{Implement（Coder：解压与信息恢复）}：Coder 把 Planner 给定的 Schema 当作不可轻易修改的设计蓝图，把其中省略的工程信息恢复出来：在若干受控 Episode 内处理语法细节、内存偏移、屏障配对和局部参数调优。它的任务和边界都很明确，不必漫无目的地探索和重试，也不再自行改变设计方向。


\item
\textbf{Optimize（Planner：沿金字塔分层演进）}：Planner 收到的是 Coder 整理后的实测反馈，不夹带编译排错的过程，注意力不会被稀释。它结合这些反馈和 \InlineCode{Microbench} 等硬件 Wiki，沿 Schema 决策金字塔自顶向下修改 Schema：先看 \InlineCode{Math\_\allowbreak{}reshape} 能否改写计算路径，再考虑 \InlineCode{Work\_\allowbreak{}schedule} 的分块映射与负载均衡，最后调整 \InlineCode{Tile\_\allowbreak{}pipeline} 中的并发掩盖，由此确定下一轮的设计方向。


\end{itemize}

\OriginalImage{assets/framework.jpg}{基于 Kernel Schema 与 AKA 的两阶段迭代框架}{\emph{图：两阶段协同优化总框架。Planner 在 Schema 层结合硬件规格、数学改写与瓶颈分析，提出并修改候选流水线（Pipeline A / B / C）；基于 AKA 的 Coder 针对选定的 Schema 生成多个实现版本；Verifier 把实测性能与诊断回传给 Planner，作为下一轮设计的依据。}}

\textbf{不同设计方案之间有继承和分叉关系。设计与代码分离后，Planner 可以在 Schema 层维护一棵“Schema 设计遗传树（Design Genetic Tree）”：}


\begin{DiagramBlock}
\IfFileExists{diagrams/diagram-10.tex}{%
\begin{adjustbox}{max width=\linewidth,max totalheight=.80\textheight,center}
\begingroup%
\def\baselinestretch{1}%
\definecolor{D10Blue}{HTML}{3478B8}%
\definecolor{D10Teal}{HTML}{27877F}%
\definecolor{D10Ink}{HTML}{243447}%
\definecolor{D10Slate}{HTML}{64748B}%
\begin{tikzpicture}[
  x=1mm, y=1mm,
  card/.style={draw, fill=white, rounded corners=2mm, line width=.75pt,
    minimum height=13mm, inner xsep=3mm, inner ysep=2.5mm, outer sep=.4pt,
    align=flush center, text=D10Ink, font=\fontsize{9.5}{12.5}\selectfont},
  entry/.style={card, draw=D10Slate, minimum width=76mm},
  schema/.style={card, draw=D10Blue, minimum width=65mm,
    font=\bfseries\fontsize{9.5}{12.5}\selectfont},
  kernel/.style={card, draw=D10Teal, minimum width=72mm},
  edge text/.style={fill=white, inner xsep=1mm, inner ysep=.8mm, align=flush center,
    text=D10Ink, font=\fontsize{9}{12.5}\selectfont},
  flow/.style={draw=D10Blue, line width=.7pt, rounded corners=1.5mm,
    -{Latex[length=2.5mm,width=1.6mm]}},
  implement/.style={flow, draw=D10Teal, dash pattern=on 2.6pt off 2pt}
]
\node[entry] (zero) at (-42,0) {从零开始：Planner · Design\\
  参考相近算子写出初始 Schema};
\node[entry] (sota) at (42,0) {从已有 SOTA 开始：Planner · Extract\\
  从现有代码提取 Schema};
\node[schema] (root) at (0,-24) {Schema v1（根节点）\\单核静态分块流水};

\node[schema] (a) at (-42,-57) {Schema v2-A（分支 A）\\2-CTA Cluster 常驻步长调度};
\node[schema] (b) at (42,-57) {Schema v2-B（分支 B）\\Split-KV 分段计算 + 归并};
\node[schema] (a1) at (-42,-92) {Schema v3-A1（子代）\\奇偶双路 Softmax 掩盖 SFU};

\node[kernel] (kb) at (42,-92) {Kernel 实现簇（v1, v2, v3…）\\
  逼近 v2-B 的实际性能上限};
\node[kernel] (ka) at (42,-139) {Kernel 实现簇（v1, v2, v3…）\\
  逼近 v2-A 的实际性能上限};
\node[kernel] (ka1) at (-42,-139) {Kernel 实现簇（v1, v2, v3…）\\
  实现双路流水的性能收益};

\coordinate (entry-turn) at (0,-12);
\coordinate (root-top-left) at ([xshift=-9mm]root.north);
\coordinate (root-top-right) at ([xshift=9mm]root.north);
\draw[flow] (zero.south) -- (zero.south |- entry-turn)
  -- (root-top-left |- entry-turn) -- (root-top-left);
\draw[flow] (sota.south) -- (sota.south |- entry-turn)
  -- (root-top-right |- entry-turn) -- (root-top-right);

\coordinate (branch-turn) at (0,-37);
\coordinate (root-bottom-left) at ([xshift=-9mm]root.south);
\coordinate (root-bottom-right) at ([xshift=9mm]root.south);
\draw[flow] (root-bottom-left) -- (root-bottom-left |- branch-turn)
  -- (a.north |- branch-turn) -- node[edge text] {Optimize：改写 Work\_schedule} (a.north);
\draw[flow] (root-bottom-right) -- (root-bottom-right |- branch-turn)
  -- (b.north |- branch-turn) -- node[edge text] {Optimize：拆分多阶段 Step} (b.north);
\draw[flow] (a.south) -- node[edge text] {Optimize：深化 Tile\_pipeline} (a1.north);

\coordinate (a-implementation) at (0,-111);
\draw[implement] (a.east) -- (a.east -| a-implementation)
  -- node[edge text, at end] {Coder · Implement} (a-implementation)
  -- (ka.west -| a-implementation) -- (ka.west);
\draw[implement] (b.south) -- node[edge text] {Coder · Implement} (kb.north);
\draw[implement] (a1.south) -- node[edge text] {Coder · Implement} (ka1.north);
\end{tikzpicture}%
\endgroup%
%
\end{adjustbox}%
}{%
\MissingDiagram{10}%
}
\end{DiagramBlock}
编译报错和指针 Bug 都在叶子节点下的 Coder 实现簇内部处理，不会污染 Planner 的上下文；\textbf{一条潜在的更高上限的设计分支，也不会因为代码暂时没优化好而被放弃}。Planner 只负责在这棵树上选择和扩展设计，Coder 只负责实现给定的节点。


\newpage

\subsection{面向长程迭代的两个机制：\InlineCode{Step} 与 \InlineCode{Dispatch}}\label{sec:step-dispatch}
把 Kernel Schema 用于复杂生产算子的长程迭代时，我们补充了两个表达机制：\textbf{Step} 描述一个算子内部的阶段划分，\textbf{Dispatch} 描述不同负载之间的设计分支。前者让 Planner 可以调整算子的阶段边界，后者让面向不同负载的特化设计不必各存一份完整副本。


\Needspace{6\baselineskip}
\subsubsection{\InlineCode{Step}：描述算子内部的多个阶段}\label{sec:step}
设计 \InlineCode{Work\_\allowbreak{}schedule} 与 \InlineCode{Tile\_\allowbreak{}pipeline} 时我们发现，工业级的复杂算子很少从头到尾是同一个循环体，\textbf{内部通常有明确的阶段边界，各阶段的调度方式也不相同}。同一个维度，在前一阶段可能是独立并行的切分轴，到下一阶段变成 Worker 内的串行递推轴，收尾时又成为跨分片的规约轴。典型例子有稀疏注意力的“稀疏索引预处理 → 主注意力计算”、长序列的“分片计算 → 全局归并”，以及融合 MoE 中依次完成路由、分块矩阵乘、激活量化与聚合的多个 \InlineCode{Step}。


在 Schema 中显式写出 \InlineCode{Step}，并规定各 \InlineCode{Step} 的启动方式和数据交接方式（片上直接传递，或经全局内存交接）之后，Schema 描述的算子不再局限于功能单一、重复循环的循环核，边界可以自由扩展。这带来两个好处。一是描述完备，多阶段的复合算子也能完整写进 Schema。二是 Planner 可以在 \InlineCode{Step} 边界上做优化：拆分一个 \InlineCode{Step}，用额外的中间结果物化换取更高的并行度；或者融合相邻的 \InlineCode{Step}，省去全局内存往返和 Kernel 启动开销。这也为以后描述和优化更复杂的 \textbf{Megakernel} 做好了准备。


\Needspace{6\baselineskip}
\subsubsection{\InlineCode{Dispatch}：不同负载的特化设计只记录差异}\label{sec:dispatch}
第~\ref{sec:challenge-delivery}~节提到，现有 Kernel Agent 倾向于堆砌分支特化来适配不同的 Workload 与 Shape。对处在不同瓶颈区间的负载采用不同设计，这个特化的思路本身是对的；难点在于如何表达这些分派。如果用简单的树结构为每个分支各存一份完整的 Schema，整体架构调整、局部机制改动和纯超参数调整会混在一起，节点之间大量重复。早期提取 MoE 时，不同层级的组合配对就产生了 50 份高度重复的副本，难以维护，也看不出分派决策的依据。


为此，我们引入结构化的 \textbf{Dispatch}：\textbf{各变体共享不变的部分，差异按改动幅度分三级记录}。


\begin{enumerate}[start=1]
\item
\textbf{整体设计重构}：只有主计算路径、工作归属或主流水结构发生根本变化时，才建立独立的 \InlineCode{<id>.\allowbreak{}yaml} 设计节点，并可通过无环的 \InlineCode{\{from: id\}} 整区复用未变化的 \InlineCode{Design} 层，避免跨文件重复；


\item
\textbf{局部机制调整}：主体结构不变、只改动局部环节（如某一步重缩放因子的共享方式）时，不派生新文件，只在 \InlineCode{dispatch.\allowbreak{}md} 中用简短的差异描述记录改动点、联动接口与适用条件；


\item
\textbf{纯超参数调整}：同一机制下的分块尺寸、Split 路数等数值调优，直接用参数条目规定遍历范围与资源预算。


\end{enumerate}

有了 \InlineCode{Dispatch}，Agent 在长程迭代中遇到某个 Shape 性能不足时，可以先判断应该做超参数遍历、局部机制调整，还是重新设计一套流水线。不同负载仍可以分别优化，共享的部分只记录一次。分派记录下来以后，我们也能观察 Agent 在什么情况下、以多大幅度做特化。Dispatch 本身是一个值得继续研究的机制。


\newpage


\subsection{工程实现与初步实验}\label{sec:planner-coder-progress}
这部分工作以工程实现为主，重点是把第\ref{sec:schema}章的 Schema 用进迭代，让整个设计 Loop 跑起来。目前这套流程已在两个代码库中实现，并完成了端到端运行：


\begin{enumerate}[start=1]
\item
\textbf{KernelSchema 独立工具链}：在 \InlineCode{KernelSchema} 仓库中实现了 \InlineCode{schema\_\allowbreak{}design}、\InlineCode{schema\_\allowbreak{}extract}、\InlineCode{schema\_\allowbreak{}optimize}、\InlineCode{schema\_\allowbreak{}implement} 四个模块，由 \InlineCode{schema\_\allowbreak{}workflow} 串成一个串行闭环，可以从零任务或已有源码出发，独立完成“设计—实现—实测—优化”循环。


\item
\textbf{AKA-Schema 分支}：在 \InlineCode{atrex-\allowbreak{}kernel-\allowbreak{}agent} 仓库中建立了集成 Schema 工作流的分支。开启 Schema 模式后，系统从原生的单层探索切换为两阶段协同：原有的多 Episode 迭代引擎改造为受控的 Implement 窗口，负责按当前 Schema 实现代码；窗口之间插入 Planner，查看反馈并更新 Schema，使 Schema 与 Kernel 一起迭代。


\end{enumerate}

实验方面，整套闭环已在真实集群上运行。受 GPU 机时限制，完整的多算子长程对照尚未展开。我们先选取一个有代表性的 MoE 算子（\InlineCode{qwen3\_\allowbreak{}30b} 配置：H2048 / I768，128 个专家、top-8，单 Token 输入），从现有 SOTA Kernel 出发，对比两阶段闭环（\InlineCode{Schema}）与原生 AKA（\InlineCode{Control}）的优化轨迹，初步效果比较明显：


\begin{itemize}
\item
\textbf{同延迟比 Token}：首次达到 ≤27.5 µs 时，\InlineCode{Schema} 组的优化阶段输出 Token 为 0.386M，\InlineCode{Control} 组为 0.745M，少用约 48\%；


\item
\textbf{同 Token 比延迟}：同为 0.4M 输出 Token 预算，\InlineCode{Schema} 组的最佳延迟为 26.70 µs，\InlineCode{Control} 组为 27.58 µs，低 3.2\%。起点已是 SOTA，延迟本身的提升空间不大。


\end{itemize}

\OriginalImage{assets/04_token_latency_combined.png}{MoE 算子：同延迟少用 Token，同 Token 降低延迟}{\emph{图：从 SOTA 出发的 MoE 算子优化，每组各一次运行。左为首次达到 ≤27.5 µs 时的输出 Token，右为 0.4M 输出 Token 预算内的最佳延迟。输出 Token 只计优化阶段已完成的 Episode（含未晋升轮次）。这是初步观察，不构成统计意义上的结论。}}


\SourceSeparator
\clearpage
\section{面向挑战三：灰盒交付、分级验证与性能预估}\label{sec:graybox}
第三个挑战在于交付与维护环节：生成的代码是黑盒，隐式缺陷难以排查，也难以确认能否用于生产。另外，第~\ref{sec:challenge-iteration}~节还留下一个问题：能否预先判断硬件极限，从而知道迭代何时可以停止。我们在这里分成三点讨论这个问题的解决思路。


\SourceSeparator
\subsection{灰盒交付：设计交给人审阅，实现交给测试把关}\label{sec:black-to-gray}
对于生成的算子难以理解和维护的问题，Kernel Schema 解决的是设计这一层：\textbf{它把原本完全黑盒的代码，变成了高层透明的“灰盒（Grey-Box）”}。底层代码仍然难读，但设计有了可读、可审的记录。


在设计层，Schema 语法紧凑，篇幅在 300～500 行，并配有事件依赖图，人类专家可以直接阅读和审查算子的算法改写与流水线决策。这些设计以结构化的形式保存下来，不再停留在 Agent 对话里的零散文字中，可以在团队内传阅、迁移到其他硬件，并在多轮迭代之间累积。长程迭代中的关键发现因此能积累成经验，供下一轮迭代和后续的算子开发使用。


在实现层，我们能保证 Schema 清晰可读，却很难要求 Implement 阶段生成的数千行代码同样便于逐行阅读。这里我们选择一条较为激进的路线：\textbf{算子最终用于真实环境，人写代码同样会犯错；底层代码既然无法靠人工逐行担保，那么在真实环境和覆盖充分的 Benchmark 中稳定可用，就代表实现是正确的。}


\SourceSeparator
\subsection{分层验证：Data/Bench 作为起点/终点}\label{sec:tiered-verification}
用验证来定义可用性，会带来新的问题。要发现 Agent 的 \textbf{Reward Hacking（钻评测漏洞）}以及共享内存竞争、越界访问等隐式 Bug，需要大规模测试集、动态边界用例和并发 Sanitizer；而实验中我们观察到，\textbf{串行等待正确性与性能验证，是拖慢迭代的最大因素}。


为兼顾测试覆盖和迭代速度，我们计划在 \InlineCode{AKA /\allowbreak{} Implement} 一侧引入类似 \textbf{Spira} 的多级门禁（Multi-Level Gates）和异步验证，并构建相应的分层测试集。


\OriginalImage{assets/spira.jpg}{Spira 多层门禁与动态评估加固示意}{\emph{图：借鉴 Spira 漏斗思想的多层门禁与动态评估加固机制。} \InlineCode{Soft Gate} \emph{快速拦截大部分低级错误，只有高潜候选进入} \InlineCode{Hard Gate} \emph{做深层验证；}\InlineCode{Advisor} \emph{可动态触发新评测以加固门禁（}\InlineCode{harden}\emph{），防范投机与隐式缺陷。}}
门禁由轻到重分为三级，通过轻量门禁的候选不必等待重型测试：


\begin{itemize}
\item
\textbf{Gate 1 快速初筛}：编译本身就是一道拦截，再加上静态规则校验和少量小 Shape 冒烟用例，以秒级代价拦截大部分语法、接口和基础逻辑错误；


\item
\textbf{Gate 2 / 3 深测与压测}：用复杂度递增的分层题目，检验动态边界上的正确性和生产负载下的综合速度（如加权 \InlineCode{Geomean}），再用重型压测和 Sanitizer 排查极端并发问题与长尾性能下降；


\item
\textbf{轻量放行、后台严测}：候选通过轻量门禁后\textbf{立即返回主循环继续探索}；后台根据候选的重要程度（如是否刷新最优记录）异步追加重型边界测试和压测，发现隐式错误或性能下降时再告警或回滚。重型测试因此不会阻塞主循环。


\end{itemize}

\begin{DiagramBlock}
\IfFileExists{diagrams/diagram-12.tex}{%
\begin{adjustbox}{max width=\linewidth,max totalheight=.80\textheight,center}
\begingroup%
\def\baselinestretch{1}%
\definecolor{D12Blue}{HTML}{3478B8}%
\definecolor{D12Teal}{HTML}{27877F}%
\definecolor{D12Amber}{HTML}{C58B28}%
\definecolor{D12Ink}{HTML}{243447}%
\definecolor{D12PaleBlue}{HTML}{F0F6FC}%
\definecolor{D12PaleTeal}{HTML}{EFF8F5}%
\begin{tikzpicture}[
  x=1mm, y=1mm,
  card/.style={draw, fill=white, rounded corners=2mm, line width=.75pt,
    inner xsep=2.5mm, inner ysep=2.5mm, outer sep=.4pt, align=flush center,
    text=D12Ink, font=\fontsize{9.5}{12.5}\selectfont},
  fast/.style={card, draw=D12Teal, minimum height=13mm},
  async/.style={card, draw=D12Amber, minimum width=72mm, minimum height=17mm},
  loop core/.style={card, draw=D12Blue, text width=42mm, minimum width=47mm, minimum height=29mm,
    font=\bfseries\fontsize{9.5}{12.5}\selectfont},
  group box/.style={draw, rounded corners=2.5mm, line width=.7pt,
    inner sep=0pt, outer sep=.4pt},
  group title/.style={inner sep=0pt, align=flush center, text=D12Ink,
    font=\bfseries\fontsize{10}{12.5}\selectfont},
  edge text/.style={fill=D12PaleTeal, inner xsep=1mm, inner ysep=.5mm, align=flush center,
    text=D12Ink, font=\fontsize{9}{12.5}\selectfont},
  flow/.style={draw=D12Teal, line width=.7pt, rounded corners=1.8mm,
    -{Latex[length=2.5mm,width=1.6mm]}},
  async flow/.style={flow, draw=D12Amber, dash pattern=on 2.6pt off 2pt},
  loop flow/.style={flow, draw=D12Blue, line width=.9pt,
    -{Latex[length=3mm,width=2mm]}}
]
\node[group title] at (0,12)
  {分级异步验证：减少串行等待，防范 Reward Hacking 与隐式 Bug};
\node[fast, minimum width=51mm, minimum height=9mm] (cand) at (-41,0)
  {Coder 生成候选 Kernel};
\node[fast, minimum width=59mm] (g1) at (-41,-21)
  {Gate 1：轻量同步检查\\
   编译、静态校验、小 Shape 测试};
\node[fast, minimum width=56mm] (next) at (45,-21)
  {返回主循环，进入下一轮\\
   （主链路不阻塞）};
\node[async] (g2) at (-41,-49)
  {Gate 2：动态边界与性能\\
   动态 Shape、掩码 / 空块边界、Geomean};
\node[async] (g3) at (-41,-75)
  {Gate 3：上线前大量的压测\\
   全量盲测、Sanitizer、掉点监控};

\node[group title] at (0,-105)
  {系统的头与尾：头部 Data 驱动探索，尾部 Bench 验证结果};
\node[loop core] (head) at (-53,-129)
  {头部：Data\\
   Kernel Schema\\与 schema\_wiki\\
   提供设计先验};
\node[loop core] (mid) at (0,-129)
  {中部：Planner–Code\\
   设计空间树搜索\\
   与受控代码实现};
\node[loop core] (tail) at (53,-129)
  {尾部：Bench\\
   分层的正确性\\与速度基准\\
   以实测结果为准};

\begin{scope}[on background layer]
  \node[group box, draw=D12Teal, fill=D12PaleTeal, fit={(-81,19) (81,-89)}] {};
  \node[group box, draw=D12Blue, fill=D12PaleBlue, fit={(-81,-98) (81,-150)}] {};
\end{scope}

\draw[flow] (cand) -- (g1);
\draw[flow] (g1) -- node[edge text] {通过即放行} (next);
\draw[async flow] (g1) -- node[edge text] {后台按优先级异步触发} (g2);
\draw[async flow] (g2) -- (g3);
\draw[async flow] (g3.east) -- (next.south |- g3.east)
  -- node[edge text] {发现缺陷则告警回滚} (next.south);

\draw[loop flow] (head) -- (mid);
\draw[loop flow] (mid) -- (tail);
\end{tikzpicture}%
\endgroup%
%
\end{adjustbox}%
}{%
\MissingDiagram{12}%
}
\end{DiagramBlock}
门禁是否可靠，取决于测试集本身。因此，我们接下来的一项重点工作，是设计一套覆盖面广、在正确性判定和速度评估（如 \InlineCode{Geomean} 与分档阈值）上分阶段、分层级组织的测试集。


从整个系统看，第\ref{sec:schema}章的 \InlineCode{schema\_\allowbreak{}wiki} 是系统的“头”，提供高质量的设计数据（\InlineCode{Data}）；本节的分层门禁与评测基准是系统的“尾”，提供高可信的验证（\InlineCode{Bench}）。\textbf{任何 Loop 都离不开这两端：以 Data 驱动探索（Data-Driven），以 Bench 验证结果（Bench-Verified）。}


\SourceSeparator
\subsection{性能预估：写代码之前判断设计的优劣}\label{sec:perf-prediction}
还有两个问题没有回答：当前结果离硬件极限（SOL，Speed-of-Light）还有多远；能否在编码和实测之前判断设计的优劣。我们计划研究的一条路径是：\textbf{基于} \InlineCode{Kernel Schema + Microbench} \textbf{构建轻量的先验性能估算器}。


\begin{DiagramBlock}
\IfFileExists{diagrams/diagram-13.tex}{%
\begin{adjustbox}{max width=\linewidth,max totalheight=.80\textheight,center}
\begingroup%
\def\baselinestretch{1}%
\definecolor{D13Blue}{HTML}{3478B8}%
\definecolor{D13Teal}{HTML}{27877F}%
\definecolor{D13Ink}{HTML}{243447}%
\definecolor{D13Slate}{HTML}{64748B}%
\begin{tikzpicture}[
  x=1mm, y=1mm,
  card/.style={draw, fill=white, rounded corners=2mm, line width=.75pt,
    minimum width=77mm, minimum height=17mm, inner xsep=2.5mm, inner ysep=2.5mm,
    outer sep=.4pt, align=flush center, text=D13Ink, font=\fontsize{9.5}{12.5}\selectfont},
  flow/.style={draw=D13Ink, line width=.7pt, rounded corners=1.5mm,
    -{Latex[length=2.5mm,width=1.6mm]}}
]
\node[card, draw=D13Slate] (in1) at (-42,0)
  {尚未实现的候选设计\\
   Schema A / B / C\\
   （数学账本、调度波次、微流水事件图）};
\node[card, draw=D13Slate] (in2) at (42,0)
  {目标硬件 Microbench 先验\\
   指令周期、吞吐、并发代价};
\node[card, draw=D13Blue, minimum height=14mm,
  font=\bfseries\fontsize{9.5}{12.5}\selectfont] (model) at (0,-33)
  {Schema + Microbench\\
   轻量先验性能模型};

\node[card, draw=D13Teal, minimum height=24mm] (abs) at (-42,-70)
  {绝对上界：估算 SOL\\
   评估可达峰值比例\\
   辅助判断平台期与停止时机};
\node[card, draw=D13Teal, minimum height=24mm] (rel) at (42,-70)
  {相对比较：近似保序\\
   不求绝对精度，只需保持优劣次序\\
   Est(A) \textless{} Est(B)\space→\space Real(A) \textless{} Real(B)};
\node[card, draw=D13Teal] (prune) at (42,-105)
  {编码前剪枝\\
   淘汰较差的方案，\\
   只把选中的 Schema 交给 Coder};

\coordinate (input-turn) at (0,-17);
\coordinate (model-top-left) at ([xshift=-17mm]model.north);
\coordinate (model-top-right) at ([xshift=17mm]model.north);
\draw[flow] (in1.south) -- (in1.south |- input-turn)
  -- (model-top-left |- input-turn) -- (model-top-left);
\draw[flow] (in2.south) -- (in2.south |- input-turn)
  -- (model-top-right |- input-turn) -- (model-top-right);

\coordinate (output-turn) at (0,-49);
\coordinate (model-bottom-left) at ([xshift=-17mm]model.south);
\coordinate (model-bottom-right) at ([xshift=17mm]model.south);
\draw[flow] (model-bottom-left) -- (model-bottom-left |- output-turn)
  -- (abs.north |- output-turn) -- (abs.north);
\draw[flow] (model-bottom-right) -- (model-bottom-right |- output-turn)
  -- (rel.north |- output-turn) -- (rel.north);
\draw[flow] (rel) -- (prune);
\end{tikzpicture}%
\endgroup%
%
\end{adjustbox}%
}{%
\MissingDiagram{13}%
}
\end{DiagramBlock}
Schema 中已经显式记录了计算与访存账本、Worker 并行映射和片内事件依赖图，再结合目标卡 \InlineCode{Microbench} 实测的指令周期与吞吐，估算器可以在两方面帮助迭代：


\begin{enumerate}[start=1]
\item
\textbf{估计绝对上界}：在给定的硬件参数和 Schema 设计下，推算由关键路径和最忙单元决定的性能上界，估计该设计能发挥多少硬件峰值，为“离物理极限还有多远、何时停止追加预算”提供一把先验标尺。


\item
\textbf{比较相对优劣}：寄存器分配、指令发射争用和缓存行为都会影响实际延迟，仿真器很难同时做到轻量、快速，又能准确预测绝对延迟。但 Planner 在尚未实现的 Schema 之间做筛选时，并不需要准确的绝对数值，\textbf{只需要近似保序（Order-Preserving）}：即使估计有系统性偏差，只要模型判定 Schema A 优于 Schema B 时，两者实现后的实测表现大概率保持同样的次序，Planner 就能在写代码之前淘汰较差的方案，保留有潜力的设计。


\end{enumerate}

\SourceSeparator
\clearpage
\section{总结与展望：从算子走向 Infra 自进化}\label{sec:outlook}
第\ref{sec:motivation}章提出了两个方向：面向当下，降低 MaaS 推理的算子成本；面向未来，作为自我改进闭环的起点。本章先小结进展，再回到第二个方向。


\subsection{工作进展小结}\label{sec:recap}
\begin{itemize}
\item
\textbf{数据与表示}（已实现）：设计语言与检查工具的压缩比在 10× 以上；首批 18 个算子、48 份设计已入库；仅在提示词中注入 Schema，Agent 就能更早达标，各 Run 结果更集中。


\item
\textbf{生成与迭代}（已实现，长程实验进行中）：Planner–Coder 闭环已端到端运行，在代表性 MoE 算子上，达到同一延迟目标少用约 48\% 的输出 Token；多算子的完整对照受机时限制尚未展开。


\item
\textbf{交付与维护}（规划中）：灰盒交付、分级异步验证，以及基于 \InlineCode{Schema + Microbench} 的性能预估。


\end{itemize}

\textbf{三部分做的是同一件事：把专家的设计直觉写成可检查、可迭代的 Schema，让 Agent 在压缩态的设计上搜索，在解压态的代码上实测。}


\SourceSeparator
\subsection{Kernel Schema 的直接效果}\label{sec:near-term}
\begin{enumerate}[start=1]
\item
\textbf{异构卡适配：改设计，不重写代码。}数学改写大多可以沿用，主要重做工作调度与微流水线。


\item
\textbf{新模型适配：从相近的设计出发。}在 \InlineCode{schema\_\allowbreak{}wiki} 中检索结构相近的算子，在它们的设计上改写。


\item
\textbf{Megakernel：融合边界成为设计决策。}\InlineCode{Step} 已能描述阶段间的切分、融合与交接，是 Schema 走向框架层的接口。


\item
\textbf{反哺训练：Schema 是蒸馏过的语料。}设计与配对源码可直接用作后训练数据，回应第~\ref{sec:challenge-data}~节的语料问题。


\end{enumerate}

\SourceSeparator
\subsection{对长程 Infra Agent 的启发}\label{sec:lessons}
放到更一般的长程 Agent 任务中，核心矛盾是：\textbf{任务的信息量，远大于 Agent 一次能有效关注的信息量。}超长程迭代和多智能体协同的困难大多源于此：上下文有限，注意力被无关信息稀释，状态难以在多轮、多角色之间可靠传递。Kernel Schema 给出两方面启发：


\begin{enumerate}[start=1]
\item
\textbf{分层抽象：每个角色只关注本层的决策。}固定的是层间交接物，层内仍是端到端搜索，并不退回 rule-based。切口应放在信息能大幅压缩、两侧搜索景观明显不同的地方（第~\ref{sec:design-space-search}~节）；否则多出来的层级只会变成 Loop Engineering，交接越多，损耗越大。


\item
\textbf{语言压缩：分层需要面向 Agent 的表示，表达即压缩。}CuTe 这类 DSL 是写给人的，复用的是底层机制；面向 Agent 的抽象应按语义分级，把下层当作黑箱，正如 CUDA C++ 封装了 PTX。设计这门语言的方法同样可以迁移：数据驱动归纳、删除测试、人机结对。


\end{enumerate}

\SourceSeparator
\subsection{展望：把方法推广到更上层的 Infra}\label{sec:infra-rsi}
Kernel 是 Infra 自我改进阶梯的第一级，也是最容易验证的一级：验证器客观，压缩比大，有 \InlineCode{Microbench} 这样的成本先验，还有大量可蒸馏的 SOTA 算子。每往上一级，设计空间更大，反馈更慢，验证更难。


\begin{DiagramBlock}
\IfFileExists{diagrams/diagram-15.tex}{%
\begin{adjustbox}{max width=\linewidth,max totalheight=.80\textheight,center}
\begingroup%
\def\baselinestretch{1}%
\definecolor{D15Teal}{HTML}{27877F}%
\definecolor{D15Amber}{HTML}{B7791F}%
\definecolor{D15Ink}{HTML}{243447}%
\definecolor{D15Slate}{HTML}{64748B}%
\definecolor{D15PaleTeal}{HTML}{EFF8F5}%
\def\Gap#1{\textcolor{D15Amber}{#1}}%
\begin{tikzpicture}[
  x=1mm, y=1mm,
  rung/.style={draw=D15Slate!60, fill=white, rounded corners=1.8mm,
    line width=.6pt, minimum width=152mm, minimum height=10.5mm, outer sep=0pt},
  name/.style={text=D15Ink, align=center, text width=22mm,
    font=\bfseries\fontsize{9.5}{11.5}\selectfont},
  cell/.style={text=D15Ink, align=center, font=\fontsize{8.5}{10.5}\selectfont},
  head/.style={text=D15Ink, align=center,
    font=\bfseries\fontsize{9}{11}\selectfont},
  note/.style={text=D15Slate, font=\fontsize{8.5}{10.5}\selectfont},
  flow/.style={-{Stealth[length=2.25mm,width=1.65mm]}, line width=.7pt}
]
  \def\cA{12}\def\cB{48}\def\cC{93}\def\cD{134}%
  \node[head] at (\cA,61) {层级};
  \node[head] at (\cB,61) {设计层\\{\mdseries（Schema 的对应物）}};
  \node[head] at (\cC,61) {验证器\\{\mdseries（单次反馈周期）}};
  \node[head] at (\cD,61) {成本先验\\{\mdseries（Microbench 的对应物）}};

  \node[rung, draw=D15Teal, fill=D15PaleTeal, line width=.9pt] at (76,0) {};
  \node[name] at (\cA,0) {Kernel\\[-.3mm]{\mdseries\fontsize{7.5}{9}\selectfont\color{D15Teal}本文起点}};
  \node[cell, text width=44mm] at (\cB,0) {数学改写 · 分块映射 · 流水编排};
  \node[cell, text width=46mm] at (\cC,0) {编译 + 数值对齐 + 实测延迟\\（秒～分钟）};
  \node[cell, text width=34mm] at (\cD,0) {Microbench\\Roofline};

  \node[rung] at (76,12) {};
  \node[name] at (\cA,12) {Megakernel};
  \node[cell, text width=44mm] at (\cB,12) {融合边界 · 阶段交接\\跨算子调度};
  \node[cell, text width=46mm] at (\cC,12) {端到端数值对齐 + 单步延迟\\（分钟）};
  \node[cell, text width=34mm] at (\cD,12) {Microbench\\\Gap{+ 跨算子干扰（待建）}};

  \node[rung] at (76,24) {};
  \node[name] at (\cA,24) {推理框架};
  \node[cell, text width=44mm] at (\cB,24) {并行策略 · 批调度\\PD 分离 · KV Cache 管理};
  \node[cell, text width=46mm] at (\cC,24) {与参考输出一致 + 吞吐 / SLO\\（分钟～小时）};
  \node[cell, text width=34mm] at (\cD,24) {计算与通信 Roofline\\\Gap{真实负载模型（待建）}};

  \node[rung] at (76,36) {};
  \node[name] at (\cA,36) {训练框架};
  \node[cell, text width=44mm] at (\cB,36) {并行切分 · 重计算\\通信重叠 · 精度选择};
  \node[cell, text width=46mm] at (\cC,36) {数值与 Loss 曲线对齐 + MFU\\（小时～天）};
  \node[cell, text width=34mm] at (\cD,36) {\Gap{保序的小规模}\\\Gap{代理实验（待建）}};

  \node[rung] at (76,48) {};
  \node[name] at (\cA,48) {自动化科研};
  \node[cell, text width=44mm] at (\cB,48) {研究假设 · 实验设计};
  \node[cell, text width=46mm] at (\cC,48) {\Gap{可复现的客观评测（多数尚缺）}\\（天～月）};
  \node[cell, text width=34mm] at (\cD,48) {\Gap{尚缺}};

  \draw[flow, D15Amber] (-5,-5) -- (-5,53);
  \node[note, text=D15Ink, rotate=90] at (-9,24) {越往上，设计空间越大，反馈越慢};

  \node[note, anchor=west] at (0,-9.5) {\Gap{琥珀色}：迁移到该层级时尚需建立的部分};
  \node[note, anchor=east] at (152,-9.5) {实现层分别对应 CUDA / CuTe 代码、框架代码与配置、实验代码};
\end{tikzpicture}%
\endgroup%
%
\end{adjustbox}%
}{%
\MissingDiagram{15}%
}
\end{DiagramBlock}
推理与训练框架都有 vLLM、SGLang、Megatron 等大量开源产物可供蒸馏，可以沿用同一条路径：找切口，定表示，配上验证与先验。


\begin{itemize}
\item
\textbf{推理框架}：同样存在邻域错位，MoE 层从 TP 改为 EP，在设计上是一步，在代码里牵动通信、权重布局和调度；难点在于真实负载与 SLO 下的评估。


\item
\textbf{训练框架}：数值、Loss 曲线和 MFU 都可测，难点在于反馈以小时乃至天计，关键是找到保序的小规模代理实验。


\item
\textbf{自动化科研}：假设、实验代码与评测可以一一对应，但验证器和成本先验大多尚缺，第一步是把它们建起来。


\end{itemize}

